% Generated mechanically from frozen input inputs/p01/ja/fields.tex.
% Do not edit this generated file; edit the bound locale source instead.

% OK, start here.
%

\setcounter{chapter}{8}
\chapter{体}



\phantomsection
\label{fields-section-phantom}



\section{序論}
\label{fields-section-introduction}

\noindent
本章では体の理論を論じる。{\it 体}とは、零でない元がすべて可逆である
非零環であったことを思い出そう。任意の非零元が単元であるから、同値な
言い方をすれば、体のイデアルは $(0)$ と $(1)$ の二つだけである。
したがって体は、後に展開する理論の多くにおいて最も単純な場合となる。

\medskip\noindent
体拡大の理論は、例えば体の射がすべて単射であるという点で、通常の可換
代数とは趣を異にする。それでも、環に関する問題はしばしば体に関する問題へ
帰着できることが分かる。例えば、任意の整域は体（その分数体）に埋め込める。
また、任意の {\it 局所環}（すなわち極大イデアルをただ一つもつ環。この用語は
まだ定義していない）には、その剰余体（すなわち極大イデアルによる商）が
付随する。したがって体拡大の知識は有用となる。




\section{基本的定義}
\label{fields-section-definitions}

\noindent
本章を可換代数を扱う章より前に置いたため、代数の諸章でさらに詳しく論じる
基本的定義のいくつかを、ここで先に導入する必要がある。

\begin{definition}
\label{fields-definition-field}
{\it 体}とは、零でない元がすべて可逆である非零環をいう。体が与えられたとき、
それ自身が体となる部分環を {\it 部分体} という。
\end{definition}

\noindent
体 $k$ に対し、部分集合 $k \setminus \{0\}$ を $k^*$ と書く。これは、環 $R$ の
可逆元の群を表す通常の記法 $R^*$ を一般化したものである。

\begin{definition}
\label{fields-definition-domain}
{\it 整域}とは、$0$ だけが零因子である非零環をいう。
\end{definition}



\section{体の例}
\label{fields-section-examples}

\noindent
まず、体の例をいくつか挙げることから始めよう。$R$ が環で $I \subset R$ が
イデアルならば、$R/I$ が体であることと $I$ が極大イデアルであることが
同値であったことを思い出してほしい。

\begin{example}[有理数]
\label{fields-example-rational-numbers}
有理数は体をなす。これは {\it 有理数体} と呼ばれ、$\mathbf{Q}$ と表される。
\end{example}

\begin{example}[素体]
\label{fields-example-prime-field}
$p$ が素数ならば、$\mathbf{Z}/(p)$ は体であり、$\mathbf{F}_p$ と表される。
実際、$(p)$ は $\mathbf{Z}$ の極大イデアルである。したがって体は有限にも
なり得る。$\mathbf{F}_p$ は $p$ 個の元をもつ。
\end{example}

\begin{example}
% P01-E0235
\label{fields-example-quotient-polymial-ring}
主イデアル整域では、既約元が生成するイデアルは極大である。さて、$k$ が
体ならば、多項式環 $k[x]$ は主イデアル整域である。したがって、$P \in k[x]$
が既約多項式（すなわち、より低い次数の項の積へ分解できない非定数多項式）
ならば、$k[x]/(P)$ は体である。これは自然な仕方で $k$ のコピーを含む。
これは体を構成する非常に一般的な方法である。例えば、複素数体
$\mathbf{C}$ は $\mathbf{R}[x]/(x^2 + 1)$ として構成できる。
\end{example}

\begin{example}[分数体]
\label{fields-example-quotient-field}
整域 $A$ が与えられると、$\mathbf{Q}$ を $\mathbf{Z}$ から構成するのとまったく
同じ方法で $A$ から構成される体 $F$ への埋め込み $A \to F$ が存在したことを
思い出そう。形式的には、$F$ の元は分数 $a/b$ の同値類である。ここで
$a, b \in A$ かつ $b \not = 0$ である。通常どおり、$a/b = a'/b'$ であることと
$ab' = ba'$ であることは同値である。体 $F$ を $A$ の {\it 商体}、
{\it 分数体}、または {\it 分数の体} という。
分数体は次の普遍性をもつ。体 $K$ への単射な環写像
$\varphi : A \to K$ が与えられると、図式
$$
\xymatrix{
F \ar[r]_\psi & K \\
A \ar[u] \ar[ru]_\varphi
}
$$
を可換にする写像 $\psi : F \to K$ が一意に存在する。実際、そのような写像の
定義は明らかである。$\psi(a/b) = \varphi(a)\varphi(b)^{-1}$ とおけばよい。
$\varphi$ の単射性により、$b \not = 0$ ならば $\varphi(b) \not = 0$ である。
\end{example}

\begin{example}[有理関数体]
\label{fields-example-field-of-rational-functions}
$k$ が体ならば、$k$ 上の有理関数体 $k(x)$ を考えられる。これは多項式環
$k[x]$ の分数体である。言い換えれば、$F, G \in k[x]$、$G \not = 0$ に対する
商 $F/G$ を、明らかな同値関係で割った集合である。
\end{example}

\begin{example}
\label{fields-example-field-of-meromorphic-functions}
$X$ をリーマン面とする。$X$ 上の有理型関数の集合を $\mathbf{C}(X)$ と表す。
関数の乗法と加法により、$\mathbf{C}(X)$ は環となる。実際には
$\mathbf{C}(X)$ が体になることが分かる。すなわち、零でない関数 $f(z)$ が
有理型ならば $1/f(z)$ も有理型である。例えば $S^2$ をリーマン球面とすると、
複素解析から、有理型関数の環 $\mathbf{C}(S^2)$ は有理関数体
$\mathbf{C}(z)$ であることが分かる。
\end{example}



\section{ベクトル空間}
\label{fields-section-vector-spaces}

\noindent
体が扱いやすい理由の一つは、体上の加群、すなわちベクトル空間の理論が
非常に単純なことである。

\begin{lemma}
\label{fields-lemma-vector-space-is-free}
$k$ が体ならば、任意の $k$-加群は自由である。
\end{lemma}

\begin{proof}
実際、線形代数により、$k$-加群（すなわちベクトル空間）$V$ は基底
$\mathcal{B} \subset V$ をもつ。この基底は、$\mathcal{B}$ 上の自由ベクトル空間
から $V$ への同型を定める。
\end{proof}

\begin{lemma}
\label{fields-lemma-field-semi-simple}
体上の加群の任意の完全列は分裂する。
\end{lemma}

\begin{proof}
任意のベクトル空間は射影加群であるから、これは補題
\ref{fields-lemma-vector-space-is-free} から従う。
\end{proof}

\noindent
これは、加群がこのように単純に振る舞うため、後の諸章で展開する理論の多くが
体についてあまり多くを述べないもう一つの理由である。補題
\ref{fields-lemma-field-semi-simple} は、$k$ を体としたときの $k$-加群の {\it 圏} に
関する主張であることに注意せよ。完全性の概念は本質的に射の理論に関わり、
すなわち純粋に圏論的な概念を用いるからである。実際これは、いわゆる
{\it アーベル圏} の内部で定式化できる。

\medskip\noindent
以後、体上の加群の研究は線形代数であり、体のイデアル論もそれほど興味深く
ないため、本章の本題である体の {\it 拡大} を研究する。


\section{体の標数}
\label{fields-section-more-fields}

\noindent
環の圏には {\it 始対象} $\mathbf{Z}$ がある。任意の環 $R$ に対し、
$\mathbf{Z}$ からそこへの写像はただ一通り存在する。体については、そのような
始対象は存在しない。
% P01-E0236
それでも、各体がそのうちちょうど一つからちょうど一通りの写像を受け取り、
他のものからは写像を受け取らないような対象の族が存在する。

\medskip\noindent
$F$ を体とする。$F$ を環とみなして、環写像
$f : \mathbf{Z} \to F$ を得る。この環写像の像は（体の部分環として）整域で
あるから、$f$ の核は $\mathbf{Z}$ の素イデアルである。したがって $f$ の核は
$(0)$、またはある素数 $p$ に対する $(p)$ のいずれかである。

\medskip\noindent
第一の場合、$f$ は単射であり、$\mathbf{Z}$ を $F$ の部分環とみなす。さらに、
$F$ の零でない元はすべて可逆なので、$p, q \in \mathbf{Z}$、$q \not = 0$ に
対して $p/q \in F$ と述べることができる。したがってこの場合、
$\mathbf{Q}$ を $F$ の部分環とみなしてよく、実際そのようにみなす。この場合、
これが $F$ の最小の部分体であることは容易に分かる。

\medskip\noindent
第二の場合、すなわち $\Ker(f) = (p)$ のとき、
$\mathbf{Z}/(p) = \mathbf{F}_p$ は $F$ の部分環である。明らかに、これは
$F$ の最小の部分体である。

\medskip\noindent
% P01-E0237
この議論により、任意の体は最小の部分体を含み、それは $\mathbf{Q}$、または
ある素数 $p$ に対する有限体 $\mathbf{F}_p$ のいずれかであることが分かる。

\begin{definition}
\label{fields-definition-characteristic}
体 $F$ の {\it 標数} は、$\mathbf{Z} \subset F$ ならば $0$ であり、
$F$ において $p = 0$ となるならば素数 $p$ である。$F$ の {\it 素体} とは
$F$ の最小の部分体をいい、標数が零ならば $\mathbf{Q} \subset F$、標数が
$p > 0$ ならば $\mathbf{F}_p \subset F$ である。
\end{definition}

\noindent
$E \subset F$ が部分体ならば、$E$ の標数が $F$ の標数と同じであることは
容易に分かる。

\begin{example}
\label{fields-example-characteristic}
$\mathbf{F}_p$ の標数は $p$ であり、$\mathbf{Q}$ の標数は $0$ である。
\end{example}


\section{体拡大}
\label{fields-section-extensions}

\noindent
一般には、体そのものよりも体の {\it 拡大} に関心がある。これは、環そのもの
ではなく固定した環上の {\it 多元環} を研究することに似ているだろう。
体の場合に都合がよいのは、次の結果により「別の体上の体」という概念が
ちょうど体拡大の概念に戻ることである。

\begin{lemma}
\label{fields-lemma-field-maps-injective}
$F$ が体で $R$ が非零環ならば、任意の環準同型
$\varphi : F \to R$ は単射である。
\end{lemma}

\begin{proof}
$a \in \Ker(\varphi)$ を零でない元とする。このとき
$\varphi(1) = \varphi(a^{-1} a) = \varphi(a^{-1}) \varphi(a) = 0$ である。
したがって $1 = \varphi(1) = 0$ となり、$R$ は零環である。
\end{proof}

\begin{definition}
\label{fields-definition-extension}
体 $F$ が体 $E$ に含まれるとき、$E$ を $F$ の {\it 拡大体} という。
$E$ が $F$ の拡大であることを $E/F$ と書く。
\end{definition}

\noindent
したがって、$F, F'$ が体で $F \to F'$ が任意の環準同型ならば、補題
\ref{fields-lemma-field-maps-injective} によりこれは単射である。そこで、少し用語を
濫用して $F'$ を $F$ の拡大体とみなせる。別の言い方をすれば、$F$ の体拡大は、
たまたま体でもある $F$-多元環にほかならない。一般の環では環準同型が必ずしも
単射でないため、状況はまったく異なる。

\medskip\noindent
$k$ を体とする。$k$ の体拡大の {\it 圏} がある。この圏の対象は拡大 $E/k$、
すなわち体の（必然的に単射な）射
$$
k \to E,
$$
である。一方、拡大 $E/k$ と $E'/k$ の間の射は $k$-多元環準同型
$E \to E'$ である。別の言い方をすれば、それは可換図式
$$
\xymatrix{
E \ar[rr] & & E' \\
& k \ar[ru] \ar[lu] &
}
$$
である。
% P01-E0238
$k$ の拡大の圏における $E \to E'$ の射の集合を $\Mor_k(E, E')$ と表す。

\begin{definition}
\label{fields-definition-tower}
体の {\it 塔} $E_n/E_{n - 1}/\ldots/E_0$ とは、体拡大の列
$E_n/E_{n - 1}$、$E_{n - 1}/E_{n - 2}$、$\ldots$、$E_1/E_0$ をいう。
\end{definition}

\noindent
体拡大の例をいくつか挙げよう。

\begin{example}
\label{fields-example-monogenic-extension}
$k$ を体とし、$P \in k[x]$ を既約多項式とする。$k[x]/(P)$ が体であることは
すでに見た（例 \ref{fields-example-quotient-polymial-ring}）。また、これは明らかな
仕方で $k$-多元環でもあるから、$k$ の拡大体である。
\end{example}

\begin{example}
\label{fields-example-field-of-meromorphic-functions-extension-C}
$X$ がリーマン面ならば、有理型関数体 $\mathbf{C}(X)$
（例 \ref{fields-example-field-of-meromorphic-functions}）は $\mathbf{C}$ の拡大体である。
実際、$\mathbf{C}$ の任意の元は、$X$ 上の有理型な、さらに正則な定数関数を
定める。
\end{example}

\noindent
$F/k$ を体拡大とする。$S \subset F$ を任意の部分集合とする。このとき、$S$ を
含む $F$ の {\it 最小} の部分拡大、すなわち $k$ を含む $F$ の部分体が存在する。
これを見るには、$S$ と $k$ を含む $F$ の部分体すべての族を考え、その共通部分を
取ればよい。これが体であることは直接確かめられる。実際、標準的な議論により、
この体は $k$ と $S$ の元を用いた有限回の初等的な代数演算、すなわち加法、乗法、
減法、除法によって得られる $F$ の元全体の集合であることが分かる。

\begin{definition}
\label{fields-definition-generated-by}
$k$ を体とする。$F/k$ が体拡大で $S \subset F$ ならば、$k$ と $S$ を含む
$F$ の最小の部分体を $k(S)$ と書く。$S$ は体拡大 $k(S)/k$ を {\it 生成する}
という。$S = \{\alpha\}$ が一元集合ならば、$k(\{\alpha\})$ の代わりに
$k(\alpha)$ と書く。有限部分集合 $S \subset F$ で $F = k(S)$ となるものが
存在するとき、$F/k$ を {\it 有限生成体拡大} という。
\end{definition}

\noindent
例えば、$\mathbf{C}$ は $\mathbf{R}$ 上 $i$ によって生成される。

\begin{exercise}
\label{fields-exercise-C-not-countably-generated}
$\mathbf{C}$ は $\mathbf{Q}$ 上可算な生成元集合をもたないことを示せ。
\end{exercise}

\noindent
次に、一つの元で生成される拡大を分類しよう。

\begin{lemma}[単拡大の分類]
\label{fields-lemma-field-extension-generated-by-one-element}
体拡大 $F/k$ が一つの元で生成されるならば、それは $k$ 上同型を除いて、
有理関数体 $k(t)/k$、または $P \in k[t]$ を既約多項式とする拡大
$k[t]/(P)$ のいずれかである。
\end{lemma}

\noindent
体拡大の最も重要な場合の多くが一つの元で生成されることを後に見るので、
これは実際に有用である。

\begin{proof}
$F = k(\alpha)$ となる $\alpha \in F$ を取る。仮定により、そのような
$\alpha$ は存在する。不定元 $t$ を $\alpha$ に送る環の射
$$
k[t] \to F
$$
がある。その像は整域であるから、核は素イデアルである。したがって、核は
$(0)$、または $P \in k[t]$ を既約多項式とする $(P)$ のいずれかである。

\medskip\noindent
核が既約な $P \in k[t]$ に対する $(P)$ ならば、この写像は $k[t]/(P)$ を
経由し、体の射 $k[t]/(P) \to F$ を誘導する。像は $\alpha$ を含むから、
この写像が全射であり、したがって同型であることが容易に分かる。この場合、
$k[t]/(P) \simeq F$ である。

\medskip\noindent
核が自明ならば、単射 $k[t] \to F$ を得る。そこで、その分数体 $k(t)$
から $F$ への射を定義できる。$R(t), Q(t) \in k[t]$ に対する商
$R(t)/Q(t)$ を $R(\alpha)/Q(\alpha)$ に送ればよい。$k[t] \to F$ が単射である
という仮定により、$Q$ が零多項式でない限り $Q(\alpha) \neq 0$ である。
$k[t]$ の分数体は有理関数体 $k(t)$ なので、像が $\alpha$ を含む射
$k(t) \to F$ を得る。したがってこれは全射であり、ゆえに同型である。
\end{proof}

\begin{lemma}
\label{fields-lemma-common-extension-field}
$k$ を体とし、$E/k$ と $F/k$ を体拡大とする。このとき共通の体拡大
$M/k$、すなわち $k$ の拡大の射 $E \to M$ と $F \to M$ が存在するような
拡大体が存在する。
\end{lemma}

\begin{proof}
$E$ が $k$ の有限生成体拡大である場合だけを証明する。一般の場合は、
ツォルンの補題を用いる型の議論により、この場合から従う（詳細は省略する）。

\medskip\noindent
まず、$E$ が $k$ の単拡大であると仮定する。補題
\ref{fields-lemma-field-extension-generated-by-one-element} により、これは
$E = k(t)$ が有理関数体であるか、ある既約多項式 $P \in k[t]$ に対して
$E = k[t]/(P)$ であることを意味する。第一の場合は $M = F(t)$ を有理関数体と
し、明らかな写像 $E \to M$ と $F \to M$ を取る。第二の場合は、$F[t]$ に
おける $P$ の像の既約因子 $Q$ を一つ選び、$M = F[t]/(Q)$ として明らかな写像
$E \to M$ と $F \to M$ を取る。

\medskip\noindent
$E = k(\alpha_1, \ldots, \alpha_n)$ ならば、$n$ に関する帰納法により、拡大
$M/k$ と写像 $F \to M$ および
$k(\alpha_1, \ldots, \alpha_{n - 1})\allowbreak \to M$ を見いだせる。前段落で論じた
単拡大の場合により、拡大
$M'\allowbreak/k(\alpha_1, \ldots, \alpha_{n - 1})$ と写像 $M \to M'$ および
$k(\alpha_1, \ldots, \alpha_n)\allowbreak \to M'$ を見いだせる。このとき $M'$ を
$k$ の拡大とみなせばよい。
\end{proof}



\section{有限次拡大}
\label{fields-section-finite-extensions}

\noindent
$F/E$ が体拡大ならば、明らかに $F$ は $E$ 上のベクトル空間でもある
（スカラー作用は単に $F$ における乗法である）。

\begin{definition}
\label{fields-definition-degree}
$F/E$ を体拡大とする。$F$ を $E$-ベクトル空間とみたときの次元をこの拡大の
{\it 次数} といい、$[F : E]$ と表す。$[F : E] < \infty$ ならば、$F$ を
$E$ の {\it 有限次} 拡大という。
\end{definition}

\begin{example}
\label{fields-example-C-over-R}
体 $\mathbf{C}$ は、基底 $1, i$ をもつ $\mathbf{R}$ 上二次元のベクトル空間で
ある。したがって $\mathbf{C}$ は $\mathbf{R}$ の次数 2 の有限次拡大である。
\end{example}

\begin{lemma}
\label{fields-lemma-finite-goes-up}
$K/E/F$ を代数的体拡大の塔とする。$K$ が $F$ 上有限次ならば、$K$ は
$E$ 上有限次である。
\end{lemma}

\begin{proof}
定義から直ちに従う。
\end{proof}

\noindent
次の二つの例では、補題
\ref{fields-lemma-field-extension-generated-by-one-element} が与える最も重要な
特別の場合について次数を考える。

\begin{example}[有理関数体の次数]
\label{fields-example-degree-rational-function-field}
$k$ を任意の体とすると、有理関数体 $k(t)$ は有限次拡大では {\it ない}。
例えば、元の集合
$\left\{t^n, n \in \mathbf{Z}\right\}$ は $k$ 上線形独立である。

\medskip\noindent
実は、$k$ が非可算ならば、$k(t)$ は $k$-ベクトル空間として {\it 非可算}
次元である。これを示すため、元の族
$\{1/(t- \alpha), \alpha \in k\} \subset k(t)$ が $k$ 上線形独立であると主張する。
これらの間に非自明な関係があれば矛盾が生じる。例えば $\mathbf{C}$ 上で考える
ならば、$\mathbf{C}$ 上の有理型関数としての $\frac{1}{t-\alpha}$ は $\alpha$ に極をもち、
他には極をもたないことから従う。したがって、$c_i \in k^*$ で
$\alpha_i \in k$ が互いに異なるとき、任意の和
$\sum c_i \frac{1}{t - \alpha_i}$ は各 $\alpha_i$ に極をもつ。特に、これは
零にはなり得ない。

\medskip\noindent
興味深いことに、これは複素数上のヒルベルトの零点定理の短い証明を与える。
少し一般化した結果については、Algebra, Theorem
\ref{algebra-theorem-uncountable-nullstellensatz} を参照せよ。
\end{example}

\begin{lemma}
\label{fields-lemma-finite-finitely-generated}
有限次体拡大は有限生成体拡大である。逆は成り立たない。
\end{lemma}

\begin{proof}
$F/E$ を有限次体拡大とする。$\alpha_1, \ldots, \alpha_n$ を $E$ 上の
ベクトル空間としての $F$ の基底とする。このとき
$F = E(\alpha_1, \ldots, \alpha_n)$ なので、$F/E$ は有限生成体拡大である。
逆が成り立たないことは、例
\ref{fields-example-degree-rational-function-field} から従う。
\end{proof}

\begin{example}[単代数拡大の次数]
\label{fields-example-degree-simple-algebraic-extension}
例 \ref{fields-example-monogenic-extension} で論じた形の一元生成体拡大 $E/k$ を考える。
言い換えれば、$P \in k[t]$ を既約多項式として $E = k[t]/(P)$ である。このとき
次数 $[E : k]$ は多項式 $P$ の次数 $d = \deg(P)$ に等しい。実際、
\begin{equation}
\label{fields-equation-P}
P = a_d t^d + a_{d - 1} t^{d - 1} + \ldots + a_0.
\end{equation}
とし、$a_d \not = 0$ とする。このとき $k[t]/(P)$ における
$1, t, \ldots, t^{d - 1}$ の像は $k$ 上線形独立である。なぜなら、それらの
関係はどれも $P$ より真に低い次数をもち得るが、$P$ はイデアル $(P)$ における
最小次数の元だからである。

\medskip\noindent
逆に、集合 $S = \{1, t, \ldots, t^{d - 1}\}$（より正確にはそれらの像）は、
ベクトル空間として $k[t]/(P)$ を張る。実際、(\ref{fields-equation-P}) により
$a_d t^d$ は $S$ の張る空間に属する。$a_d$ は可逆なので、$t^d$ も $S$ の
張る空間に属する。同様に、関係 $t P(t) = 0$ は $t^{d + 1}$ の像が
$\{1, t, \ldots, t^d\}$ の張る空間に属することを示す。直前に示したことに
より、これは $S$ の張る空間に属する。上向きの帰納法により、$n \geq d$ に
対する $t^n$ の像は $S$ の張る空間に属することが分かる。
\end{example}

\noindent
例えば、これは $[\mathbf{C}: \mathbf{R}] = 2$ という観察を確認する。
より一般に、$k$ が体で $\alpha \in k$ が平方でないならば、既約多項式
$x^2 - \alpha \in k[x]$ により次数二の拡大 $k[x]/(x^2 - \alpha)$ を構成できる。
これを $k(\sqrt{\alpha})$ と書く。このような拡大は、明らかな理由により
{\it 二次拡大} と呼ばれる。

\medskip\noindent
次数についての基本的事実は、次数が {\it 塔において乗法的} だということである。

\begin{lemma}[乗法性]
\label{fields-lemma-multiplicativity-degrees}
体の塔 $F/E/k$ が与えられているとする。このとき
$$
[F:k] = [F:E][E:k]
$$
である。
\end{lemma}

\begin{proof}
$\alpha_1, \ldots, \alpha_n \in F$ を $F$ の $E$-基底とし、
$\beta_1, \ldots, \beta_m \in E$ を $E$ の $k$-基底とする。主張は、積の集合
$\{\alpha_i \beta_j, 1 \leq i \leq n, 1 \leq j \leq m\}$ が $F$ の
$k$-基底になるということである。まず、これらが $k$ 上 $F$ を張ることを
確かめよう。

\medskip\noindent
仮定により、$\{\alpha_i\}$ は $E$ 上 $F$ を張る。したがって $f \in F$ ならば、
ある $a_i \in E$ が存在して
$$
f = \sum\nolimits_i a_i \alpha_i,
$$
となる。また各 $i$ に対し、ある $b_{ij} \in k$ を用いて
$a_i = \sum b_{ij} \beta_j$ と書ける。これらを合わせれば
$$
f = \sum\nolimits_{i,j} b_{ij} \alpha_i \beta_j,
$$
を得るので、$\{\alpha_i \beta_j\}$ は $k$ 上 $F$ を張る。

\medskip\noindent
次に、$c_{ij} \in k$ に対する非自明な関係
$$
\sum\nolimits_{i,j} c_{ij} \alpha_i \beta_j = 0
$$
が存在したと仮定する。この場合、
$$
\sum\nolimits_i \alpha_i \left( \sum\nolimits_j c_{ij} \beta_j \right) = 0,
$$
となり、$\beta_j$ が $E$ に属するので内側の項も同じ体に属する。
$\{\alpha_i\}$ の $E$-線形独立性により、内側の和はすべて零となる。
次に $\{\beta_j\}$ の $k$-線形独立性により、$c_{ij}$ はすべて零となる。
\end{proof}

\noindent
ここで少し脇道にそれて、関連の薄い定義を一つ述べる。

\begin{definition}
\label{fields-definition-number-field}
体 $K$ の標数が $0$ で、拡大 $K/\mathbf{Q}$ が有限次であるとき、その体を
{\it 数体} という。
\end{definition}

\noindent
数体は代数的整数論の基本的対象である。後に、数体における整数
$\mathbf{Z}$ の類似物について、一意分解に似た性質がなお成り立つことを見る
（ただし通常、厳密な一意分解は成り立たない）。


\section{代数拡大}
\label{fields-section-algebraic-extensions}

\noindent
% P01-E0239
重要な拡大の一群は、すべての元が有限次拡大を生成するような拡大である。

\begin{definition}
\label{fields-definition-algebraic}
体拡大 $F/E$ を考える。元 $\alpha \in F$ が $E$ 係数のある非零多項式の根で
あるとき、$\alpha$ は $E$ 上 {\it 代数的} であるという。$F$ のすべての元が
代数的ならば、$F$ を $E$ の {\it 代数拡大} という。
\end{definition}

\noindent
補題 \ref{fields-lemma-field-extension-generated-by-one-element} により、部分拡大
$E(\alpha)$ は、有理関数体 $E(t)$、または $P \in E[t]$ を既約多項式とする
商環 $E[t]/(P)$ のいずれかと同型である。後者の場合、$\alpha$ は $E$ 上
代数的である（実際、補題
\ref{fields-lemma-field-extension-generated-by-one-element} の証明は、$\alpha$ が
$P$ の根となるように $P$ を選べることを示す）。前者の場合は代数的でない。

\begin{example}
\label{fields-example-C-algebraic-over-R}
体 $\mathbf{C}$ は $\mathbf{R}$ 上代数的である。実際、
$\mathbf{C}$ における $\alpha = a + ib$ に対し、
$\alpha^2 - 2a\alpha + a^2 + b^2 = 0$ は $\mathbf{R}$ 上の $\alpha$ に対する
多項式方程式である。
\end{example}

\begin{example}
\label{fields-example-compact-riemann-surface-is-finite-over-P1}
$X$ をコンパクトリーマン面とし、$f \in \mathbf{C}(X) - \mathbf{C}$ を $X$ 上の
任意の非定数有理型関数とする（例
\ref{fields-example-field-of-meromorphic-functions} を参照）。このとき
$\mathbf{C}(X)$ は、$f$ が生成する部分拡大 $\mathbf{C}(f)$ 上代数的であることが
知られている。ここでは証明しない。
\end{example}

\begin{lemma}
\label{fields-lemma-algebraic-goes-up}
$K/E/F$ を体拡大の塔とする。
\begin{enumerate}
\item $\alpha \in K$ が $F$ 上代数的ならば、$\alpha$ は $E$ 上代数的である。
\item $K$ が $F$ 上代数的ならば、$K$ は $E$ 上代数的である。
\end{enumerate}
\end{lemma}

\begin{proof}
定義から直ちに従う。
\end{proof}

\noindent
次に、有限性と代数性の間に深い関係があることを示す。

\begin{lemma}
\label{fields-lemma-finite-is-algebraic}
有限次拡大は代数拡大である。実際、拡大 $E/k$ が代数的であることと、ある
$\alpha \in E$ が生成する任意の部分拡大 $k(\alpha)/k$ が有限次であることは
同値である。
\end{lemma}

\noindent
一般には、代数拡大が有限次であるという主張はまったく成り立たない。

\begin{proof}
$E/k$ を次数 $n$ の有限次拡大とする。$\alpha \in E$ を選ぶ。このとき元
$1, \alpha, \ldots, \alpha^n$ は $k$ 上線形従属である。そうでなければ必然的に
$[E : k] > n$ となるからである。この線形従属関係から、$\alpha$ が満たすべき
多項式を得る。

\medskip\noindent
最後の主張について、一元生成拡大 $k(\alpha)/k$ が有限次であることと
$\alpha$ が $k$ 上代数的であることは、例
\ref{fields-example-degree-rational-function-field} および
\ref{fields-example-degree-simple-algebraic-extension} により同値である。したがって
$E/k$ が代数的ならば、各 $\alpha \in E$ に対する $k(\alpha)/k$ は有限次拡大で
あり、逆も成り立つ。
\end{proof}

\noindent
直前の証明（実際には例 \ref{fields-example-degree-rational-function-field} および
\ref{fields-example-degree-simple-algebraic-extension}）から、一元生成拡大が有限次で
あることと代数的であることが同値だという補題を取り出せる。次の結果では
この観察を用いる。

\begin{lemma}
\label{fields-lemma-algebraic-finitely-generated}
\begin{slogan}
有限生成代数拡大は有限次である。
\end{slogan}
$k$ を体とし、$\alpha_1, \alpha_2, \ldots, \alpha_n$ を、ある拡大体に属する
元で、各 $\alpha_i$ が $k$ 上代数的なものとする。このとき拡大
$k(\alpha_1, \ldots, \alpha_n)/k$ は有限次である。すなわち、有限生成代数拡大は
有限次である。
\end{lemma}

\begin{proof}
各拡大
$k(\alpha_{1}, \ldots, \alpha_{i+1})/k(\alpha_1, \ldots, \alpha_{i})$
は一つの元で生成され、かつ代数的なので有限次である。次数の乗法性
（補題 \ref{fields-lemma-multiplicativity-degrees}）により結果を得る。
\end{proof}

\noindent
% P01-E0240
$\mathbf{Q}$ 上代数的な複素数の集合を、単に {\it 代数的数} という。例えば、
$\sqrt{2}$ は代数的であり、$i$ も代数的だが、$\pi$ は代数的でない。
代数的数が体をなすことは基本的事実である。しかし、数が代数的であることを、
有理数係数の非零多項式方程式を満たすこととして定義しただけでは、これを
（例えば多項式方程式によって）証明する方法は明らかでない。

\begin{lemma}
\label{fields-lemma-algebraic-elements}
$E/k$ を体拡大とする。このとき、$E$ の元で $k$ 上代数的なものは $E/k$ の
部分拡大をなす。
\end{lemma}

\begin{proof}
$\alpha, \beta \in E$ を $k$ 上代数的とする。補題
\ref{fields-lemma-algebraic-finitely-generated} により、$k(\alpha, \beta)/k$ は
有限次拡大である。したがって
$k(\alpha + \beta) \subset k(\alpha, \beta)$ は有限次拡大であり、補題
\ref{fields-lemma-finite-is-algebraic} により $\alpha + \beta$ は代数的である。
$\alpha$ と $\beta$ の差、積、商についても同様である。
\end{proof}

\noindent
環の場合と同じく、体拡大の好ましい性質の多くは、塔と合成体によって保存される。

\begin{lemma}
\label{fields-lemma-algebraic-permanence}
$E/k$ と $F/E$ を代数的体拡大とする。このとき $F/k$ は代数的体拡大である。
\end{lemma}

\begin{proof}
$\alpha \in F$ を選ぶ。このとき $\alpha$ は $E$ 上代数的である。
% P01-E0241
重要な観察は、$\alpha$ が、$k$ のある有限生成拡大上代数的だということである。
すなわち、$\alpha$ が $k(S)$ 上代数的となるような有限集合 $S \subset E$ が
存在する。実際、代数的であるとは、$\alpha$ が満たすある多項式が $E[x]$ に
存在するということであり、その係数を $S$ として取ればよい。したがって
$\alpha$ は $k(S)$ 上代数的である。特に、拡大 $k(S, \alpha)/ k(S)$ は有限次で
ある。$S$ は有限集合で $k(S)/k$ は代数的なので、補題
\ref{fields-lemma-algebraic-finitely-generated} により $k(S)/k$ は有限次である。
乗法性（補題 \ref{fields-lemma-multiplicativity-degrees}）を用いると
$k(S,\alpha)/k$ は有限次となるので、$\alpha$ は $k$ 上代数的である。
\end{proof}

\noindent
直前の議論で用いた証明法、すなわち $E$ 上代数的であるという性質が $E$ の
有限生成部分拡大へ {\it 降下した} という考え方は、代数学全体で繰り返し現れる。
例えば、これにより一般の可換代数の問題をネーター的な場合へ帰着できることが
多い。

\begin{lemma}
\label{fields-lemma-size-algebraic-extension}
$E/F$ を代数的体拡大とする。このとき $E$ の濃度 $|E|$ は高々
$\max(\aleph_0, |F|)$ である。
\end{lemma}

\begin{proof}
$S$ を $F$ 係数の非定数多項式全体の集合とする。各 $P \in S$ に対し、根の集合
$r(P, E) = \{\alpha \in E \mid P(\alpha) = 0\}$ は有限である（詳細は省略する）。
さらに、$E$ が $F$ 上代数的であることから
$E = \bigcup_{P \in S} r(P, E)$ となる。$S$ は $F$ のコピーの有限積の可算和なので、
その濃度が $\max(\aleph_0, |F|)$ で抑えられることは明らかである。したがって
$E$ の濃度も同じ上界をもつ。
\end{proof}

\begin{lemma}
\label{fields-lemma-subalgebra-algebraic-extension-field}
$E/F$ を有限次、またはより一般に代数的な体拡大とする。任意の部分環
$F \subset R \subset E$ は体である。
\end{lemma}

\begin{proof}
$\alpha \in R$ を零でない元とする。このとき $1, \alpha, \alpha^2, \ldots$ は
$R$ に含まれる。補題 \ref{fields-lemma-finite-is-algebraic} により、$a_i \in F$ に
対する非自明な関係
$a_0 + a_1 \alpha + \ldots + a_d \alpha^d = 0$ が存在する。$a_0 \not = 0$ と
仮定してよい。そうでなければ、この関係を $\alpha$ で割って $d$ を減らせる。
すると
$$
a_0 = \alpha (- a_1  - \ldots - a_d \alpha^{d - 1})
$$
となるので、$\alpha$ の逆元は $R$ の元
$a_0^{-1} (- a_1  - \ldots - a_d \alpha^{d - 1})$ である。
\end{proof}

\begin{lemma}
\label{fields-lemma-algebraic-extension-self-map}
% P01-E0242
$E/F$ を代数的体拡大とする。任意の $F$-多元環写像
$f : E \to E$ は自己同型である。
\end{lemma}

\begin{proof}
$E/F$ が有限次ならば、$f : E \to E$ は有限次元ベクトル空間の単射な
$F$-線形写像（補題 \ref{fields-lemma-field-maps-injective}）なので、全単射である。
一般の場合でも $f$ が単射であることは分かる。$\alpha \in E$ とし、
$P(\alpha) = 0$ を満たす多項式 $P \in F[x]$ を取る。$E' \subset E$ を、
$E$ における $P$ の根 $\alpha = \alpha_1, \ldots, \alpha_n$ が生成する
$E$ の部分体とする。補題 \ref{fields-lemma-algebraic-finitely-generated} により、
$E'$ は $F$ 上有限次である。$f$ は根の集合を保存するので
$f|_{E'} : E' \to E'$ となる。したがって証明の第一部により $f|_{E'}$ は
同型であり、$\alpha$ が $f$ の像に属することが分かる。
\end{proof}






\section{最小多項式}
\label{fields-section-minimal-polynomials}

\noindent
$E/k$ を体拡大とし、$\alpha \in E$ を $k$ 上代数的とする。このとき $\alpha$ は
$k[x]$ における（非自明な）多項式方程式を満たす。$P(\alpha) = 0$ となる
多項式 $P \in k[x]$ 全体の集合を考える。仮定により、この集合は零多項式だけを
含むのではない。この集合が {\it イデアル} であることは容易に分かる。
実際、これは写像
$$
k[x] \to E, \quad x \mapsto \alpha
$$
の核である。$k[x]$ は主イデアル整域なので、このイデアルの {\it 生成元}
$P \in k[x]$ が存在する。一般性を失わず $P$ がモニックであると仮定すれば、
$P$ は一意に定まる。

\begin{definition}
\label{fields-definition-minimal-polynomial}
上の多項式 $P$ を、$k$ 上の $\alpha$ の {\it 最小多項式} という。
\end{definition}

\noindent
最小多項式は次のように特徴づけられる。これは $\alpha$ を零化する最小次数の
モニック多項式である。$P$ の任意の非定数倍はより大きな次数をもち、
$P$ の倍数だけが $\alpha$ を零化できる。これが {\it 最小} という名の理由である。

\medskip\noindent
明らかに最小多項式は {\it 既約} である。これは、$\alpha$ を零化する多項式から
なる $k[x]$ のイデアルが素イデアルだという主張と同値である。写像
$k[x] \to E, x \mapsto \alpha$ は整域（実際には体）への写像なので、その核が
素イデアルであることから従う。

\begin{lemma}
\label{fields-lemma-degree-minimal-polynomial}
最小多項式の次数は $[k(\alpha) : k]$ である。
\end{lemma}

\begin{proof}
これは補題 \ref{fields-lemma-field-extension-generated-by-one-element} の議論を
言い換えただけである。$P$ が $\alpha$ の最小多項式ならば、先の証明と同様に
写像
$$
k[x]/(P) \to k(\alpha), \quad x \mapsto \alpha
$$
は同型である。また、このような拡大の次数はすでに数えた
（例 \ref{fields-example-degree-simple-algebraic-extension} を参照）。
\end{proof}

\noindent
したがって、上の証明の観察は、$\alpha \in E$ が代数的ならば
$k(\alpha) \subset E$ は $k[x]/(P)$ と同型だということである。


\section{代数閉包}
\label{fields-section-algebraic-closure}

\noindent
「代数学の基本定理」は $\mathbf{C}$ が代数閉であると主張する。この結果の美しい
証明は複素解析におけるリウヴィルの定理を用いる。
% P01-E0243
ここでは別の証明を与える（補題 \ref{fields-lemma-C-algebraically-closed} を参照）。

\begin{definition}
\label{fields-definition-algebraically-closed}
体 $F$ の任意の代数拡大 $E/F$ が自明、すなわち $E = F$ であるとき、その体を
{\it 代数閉} という。
\end{definition}

\noindent
これはすべての文献で採用される定義とは限らない。次の補題は、これと別の
定義とを比較する。

\begin{lemma}
\label{fields-lemma-algebraically-closed}
$F$ を体とする。次の条件は同値である。
\begin{enumerate}
\item $F$ は代数閉である。
\item $F$ 上の任意の既約多項式は一次式である。
\item $F$ 上の任意の非定数多項式は根をもつ。
\item $F$ 上の任意の非定数多項式は一次因子の積である。
\end{enumerate}
\end{lemma}

\begin{proof}
$F$ が代数閉ならば、任意の既約多項式は一次式である。実際、次数 $> 1$ の
既約多項式が存在すれば、それは非自明な有限次（したがって代数的な）体拡大を
生成する。例 \ref{fields-example-degree-simple-algebraic-extension} を参照せよ。
よって (1) は (2) を含意する。任意の既約多項式が一次式ならば、任意の
既約多項式は根をもち、したがって任意の非定数多項式が根をもつ。よって
(2) は (3) を含意する。

\medskip\noindent
任意の非定数多項式が根をもつと仮定する。$P \in F[x]$ を非定数とする。
$\alpha \in F$ に対して $P(\alpha) = 0$ ならば、ある $Q \in F[x]$ により
$P = (x - \alpha)Q$ となる（余りを伴う除法による）。したがって次数に関する
帰納法により、任意の非定数多項式は積 $c \prod (x - \alpha_i)$ として書ける。

\medskip\noindent
最後に、$F$ 上の任意の非定数多項式が一次因子の積であると仮定する。
$E/F$ を代数拡大とする。このとき、$E$ のすべての単部分拡大
$F(\alpha)/F$ は必然的に自明である（仮定により既約多項式は一次式だけだから
である）。したがって $E = F$ である。(4) が (1) を含意することが分かり、
証明が完了する。
\end{proof}

\noindent
次に、体の「普遍的」な代数拡大を定義したい。実際には注意が必要である。
代数閉包は普遍対象では {\it ない}。すなわち、代数閉包は {\it 一意な} 同型を
除いて一意なのではなく、同型を除いてのみ一意である。それでも、関手的では
ないものの非常に便利である。

\begin{definition}
\label{fields-definition-algebraic-closure}
$F$ を体とする。$F$ の {\it 代数閉包} とは、$F$ を含む体 $\overline{F}$ で、
次を満たすものをいう。
\begin{enumerate}
\item $\overline{F}$ は $F$ 上代数的である。
\item $\overline{F}$ は代数閉である。
\end{enumerate}
\end{definition}

\noindent
$F$ が代数閉ならば、$F$ はそれ自身の代数閉包である。次に基本的な存在結果を
証明する。

\begin{theorem}
\label{fields-theorem-existence-algebraic-closure}
任意の体は代数閉包をもつ。
\end{theorem}

\noindent
この証明は、本章の残りの部分にとってほとんど横道である。しかし、体を代数閉体
へ埋め込むことが {\it 可能} であるとは知っておきたく、しばしばそのような
埋め込みがすでに選ばれていると仮定する。

\begin{proof}
$F$ を体とする。補題 \ref{fields-lemma-size-algebraic-extension} により、$F$ の
代数拡大の濃度は $\max(\aleph_0, |F|)$ で抑えられる。$F$ を含み、
$|S| > \max(\aleph_0, |F|)$ を満たす集合 $S$ を選ぶ。次の条件を満たす三つ組
$(E, \sigma_E, \mu_E)$ を考える。
\begin{enumerate}
\item $E$ は $F \subset E \subset S$ を満たす集合である。
\item $\sigma_E : E \times E \to E$ および $\mu_E : E \times E \to E$ は、
$(E, \sigma_E, \mu_E)$ が $F$ の体拡大の構造を定めるような集合の写像である
（特に、$a, b \in F$ に対して $\sigma_E(a, b) = a +_F b$ であり、
$\mu_E$ についても同様である）。
\item $E/F$ は代数的体拡大である。
\end{enumerate}
このような三つ組 $(E, \sigma_E, \mu_E)$ 全体は集合 $I$ をなす。
% P01-E0244
$i \in I$ に対し、対応する $F$ の体拡大を
$E_i = (E_i, \sigma_i, \mu_i)$ と表す。$i \leq i'$ であることを
$E_i \subset E_{i'}$（$E_i$ と $E_{i'}$ は同じ集合 $S$ の部分集合なので、これは
意味をもつ）、かつ
$\sigma_i = \sigma_{i'}|_{E_i \times E_i}$ および
$\mu_i = \mu_{i'}|_{E_i \times E_i}$ が成り立つこととして、$I$ に半順序を
定める。言い換えれば、$E_{i'}$ は $E_i$ の体拡大である。

\medskip\noindent
$T \subset I$ を全順序部分集合とする。このとき、誘導される写像
$\sigma_T = \bigcup \sigma_i$ および $\mu_T = \bigcup \mu_i$ を備えた
$E_T = \bigcup_{i \in T} E_i$ が $I$ の別の元になることは明らかである。
% P01-E0245
言い換えれば、$I$ の任意の全順序部分集合は $I$ に上界をもつ。ツォルンの補題に
より、$I$ には極大元 $(E, \sigma_E, \mu_E)$ が存在する。$E$ が代数閉包であると
主張する。$I$ の定義により拡大 $E/F$ は代数的なので、$E$ が代数閉であることを
示せば十分である。

\medskip\noindent
これを見るため背理法で論じる。$E$ が代数閉でないと仮定する。このとき補題
\ref{fields-lemma-algebraically-closed} により、$E$ 上に次数 $> 1$ の既約多項式 $P$ が
存在する。補題 \ref{fields-lemma-finite-is-algebraic} により、非自明な有限次拡大
$E' = E[x]/(P)$ を得る。補題 \ref{fields-lemma-algebraic-permanence} により
$E'/F$ は代数的であることに注意せよ。したがって $E'$ の濃度は
$\leq \max(\aleph_0, |F|)$ である。初等的集合論により、与えられた単射
$E \subset S$ を単射 $E' \to S$ へ拡張できる。言い換えれば、$E'$ を集合
$I$ の元とみなせるが、これは $E$ の極大性に反する。この矛盾により証明が
完了する。
\end{proof}

\begin{lemma}
\label{fields-lemma-map-into-algebraic-closure}
$F$ を体とする。$\overline{F}$ を $F$ の代数閉包とし、$M/F$ を代数拡大とする。
このとき $F$-拡大の射 $M \to \overline{F}$ が存在する。
\end{lemma}

\begin{proof}
$F \subset E \subset M$ が部分拡大で、$\varphi : E \to \overline{F}$ が
$F$-拡大の射となる対 $(E, \varphi)$ 全体の集合を $I$ とする。
$(E, \varphi) \leq (E', \varphi')$ であることを $E \subset E'$ かつ
$\varphi'|_E = \varphi$ が成り立つこととして、集合 $I$ を半順序づける。
$T = \{(E_t,  \varphi_t)\} \subset I$ が全順序部分集合ならば、
$\bigcup \varphi_t : \bigcup E_t \to \overline{F}$ は $I$ の元である。
したがって $I$ の任意の全順序部分集合は上界をもつ。ツォルンの補題により、
$I$ には極大元 $(E, \varphi)$ が存在する。$E = M$ と主張する。これを示せば
証明は完了する。そうでないならば、$\alpha \in M$、$\alpha \not \in E$ を
選ぶ。
% P01-E0246
補題 \ref{fields-lemma-algebraic-goes-up} により、$\alpha$ は $E$ 上代数的である。
$P$ を $E$ 上の $\alpha$ の最小多項式とし、$P^\varphi$ を $\varphi$ による
$P$ の $\overline{F}[x]$ における像とする。$\overline{F}$ は代数閉なので、
$P^\varphi$ は $\overline{F}$ に根 $\beta$ をもつ。そこで $x$ を $\beta$ に
送ることにより、$\varphi$ を
$\varphi' : E(\alpha) = E[x]/(P) \to \overline{F}$ へ拡張できる。
これは $(E, \varphi)$ の極大性に反し、所望の結論を得る。
\end{proof}

\begin{lemma}
\label{fields-lemma-algebraic-closures-isomorphic}
一つの体の任意の二つの代数閉包は同型である。
\end{lemma}

\begin{proof}
$F$ を体とする。$M$ と $\overline{F}$ が $F$ の代数閉包ならば、補題
\ref{fields-lemma-map-into-algebraic-closure} により、$F$-拡大の射
$\varphi : M \to \overline{F}$ が存在する。像 $\varphi(M)$ は代数閉である。
一方、補題 \ref{fields-lemma-algebraic-goes-up} により、拡大
$\varphi(M) \subset \overline{F}$ は代数的である。したがって
$\varphi(M) = \overline{F}$ である。
\end{proof}





\section{互いに素な多項式}
\label{fields-section-relatively-prime}

\noindent
$K$ を代数閉体とする。補題 \ref{fields-lemma-algebraically-closed} で見たように、環
$K[x]$ のイデアル構造は非常に単純である。特に、任意の多項式 $P \in K[x]$ は
$$
P = c(x - \alpha_1) \ldots (x - \alpha_n),
$$
と書ける。
% P01-E0247
ここで $c$ は定数項であり、
% P01-E0248
$\alpha_1, \ldots, \alpha_n \in k$ は $P$ の根である（重複度を込めて数える）。
明らかに、$K[x]$ の既約多項式は一次多項式 $c(x - \alpha)$、
$c, \alpha \in K$（かつ $c \neq 0$）だけである。

\begin{definition}
\label{fields-definition-relatively-prime}
$k$ を任意の体とする。$k[x]$ の二つの多項式が $k[x]$ の単位イデアルを生成する
とき、それらは {\it 互いに素} であるという。
\end{definition}

\noindent
上の議論を続けると、$K$ が代数閉体ならば、$K[x]$ の二つの多項式が互いに素で
あることと、それらが共通根をもたないことは同値である。実際、$K[x]$ の
極大イデアルは $(x - \alpha)$、$\alpha \in K$ の形である。したがって
$F, G \in K[x]$ が共通根をもたないならば、$(F, G)$ はどの
$(x - \alpha)$ にも含まれない（含まれれば $\alpha$ が共通根になる）。

\medskip\noindent
$k$ が代数閉で {\it ない} 場合でも、これは $k[x]$ の二つの多項式がいつ単位
イデアルを生成するかについて情報を与える。

\begin{lemma}
\label{fields-lemma-relatively-prime-polynomials}
$k[x]$ の二つの多項式が互いに素であることと、$k$ の代数閉包
$\overline{k}$ に共通根をもたないことは同値である。
\end{lemma}

\begin{proof}
主張は、任意の二つの多項式 $P, Q$ が $k[x]$ で $(1)$ を生成することと、
$\overline{k}[x]$ で $(1)$ を生成することが同値だということである。これは
線形代数の一事実である。$k$ 係数の連立一次方程式が解をもつことと、$k$ の
任意の拡大で解をもつことは同値である。したがって代数閉体の場合へ帰着でき、
その場合、結果はすでに証明したことから明らかである。
\end{proof}





\section{分離代数拡大}
\label{fields-section-separable-extensions}

\noindent
標数 $p$ では、体上の既約多項式について少し特異な現象が起こる。次の補題で
これを説明する。

\begin{lemma}
\label{fields-lemma-irreducible-polynomials}
$F$ を体とし、$P \in F[x]$ を $F$ 上の既約多項式とする。
$P' = \text{d}P/\text{d}x$ を $x$ に関する $P$ の導関数とする。このとき、
次の二つの場合のいずれかが起こる。
\begin{enumerate}
\item $P$ と $P'$ は互いに素である。または、
\item $P'$ は零多項式である。
\end{enumerate}
第二の場合が起こりうるのは、$F$ の標数が $p > 0$ のときに限る。この場合、
$q = p^f$ を $p$ の冪として $P(x) = Q(x^q)$ であり、$Q \in F[x]$ は
$Q$ と $Q'$ が互いに素となる既約多項式である。
\end{lemma}

\begin{proof}
$P'$ の次数は $< \deg(P)$ であることに注意する。したがって $P$ と $P'$ が
互いに素でないならば、$(P, P') = (R)$ であり、$R$ は次数 $< \deg(P)$ の
多項式となって $P$ の既約性に反する。これにより (1) と (2) の二者択一が
得られる。

\medskip\noindent
(2) の場合にあり、$P = a_d x^d + \ldots + a_0$ と仮定する。このとき
$P' = da_d x^{d - 1} + \ldots + a_1$ である。標数 $0$ では、これにより
$a_d, \ldots, a_1 = 0$ が強制され、$P$ が定数となるので矛盾である。
したがって標数 $p$ は正である。この場合、条件 $P' = 0$ により、$p$ が
$i$ を割り切らないときは必ず $a_i = 0$ となる。
言い換えれば、ある非定数多項式 $P_1$ により $P(x) = P_1(x^p)$ である。
明らかに $P_1$ も既約である。次数に関する帰納法により、補題の主張のように
$P_1(x) = Q(x^q)$ となることが分かる。したがって
$P(x) = Q(x^{pq})$ であり、補題が証明された。
\end{proof}

\begin{definition}
\label{fields-definition-separable}
$F$ を体とし、$K/F$ を体拡大とする。
\begin{enumerate}
\item $F$ 上の既約多項式 $P$ がその導関数と互いに素であるとき、その多項式は
{\it 分離的} であるという。
\item $F$ 上代数的な $\alpha \in K$ が与えられたとき、その最小多項式が
$F$ 上分離的ならば、$\alpha$ は $F$ 上 {\it 分離的} であるという。
\item $K$ が $F$ の代数拡大であるとき、$K$ のすべての元が $F$ 上分離的ならば、
$K$ は $F$ 上 {\it 分離的}\footnote{代数的でない拡大については、この定義は
意味をなさず、正しい定義でもない。読者には Algebra の
Sections \ref{algebra-section-separability} および
\ref{algebra-section-separability-continued} を参照されたい。}
であるという。
\end{enumerate}
\end{definition}

\noindent
補題 \ref{fields-lemma-irreducible-polynomials} により、標数 $0$ ではすべての
既約多項式が分離的であり、拡大のすべての代数的元が分離的であり、すべての
代数拡大が分離的である。

\begin{lemma}
\label{fields-lemma-separable-goes-up}
$K/E/F$ を代数的体拡大の塔とする。
\begin{enumerate}
\item $\alpha \in K$ が $F$ 上分離的ならば、$\alpha$ は $E$ 上分離的である。
\item $K$ が $F$ 上分離的ならば、$K$ は $E$ 上分離的である。
\end{enumerate}
\end{lemma}

\begin{proof}
補題 \ref{fields-lemma-irreducible-polynomials} を以下では断りなく用いる。
$P$ を $F$ 上の $\alpha$ の最小多項式とし、$Q$ を $E$ 上の $\alpha$ の
最小多項式とする。このとき $Q$ は多項式環 $E[x]$ で $P$ を割り、
$P = QR$ と書ける。すると $P' = Q'R + QR'$ である。したがって
$Q' = 0$ ならば $Q$ は $P$ と $P'$ の両方を割り、補題により $P' = 0$ となる。
これで (1) が証明された。(2) は (1) と定義から直ちに従う。
\end{proof}

\begin{lemma}
\label{fields-lemma-recognize-separable}
$F$ を体とする。$F$ 上の既約多項式 $P$ が分離的であることと、$F$ の
代数閉包において $P$ の根が相異なることは同値である。
\end{lemma}

\begin{proof}
$\alpha \in \overline{F}$ が $P$ と $P'$ の両方の根であると仮定する。
このとき、ある多項式 $Q$ により $P = (x - \alpha)Q$ である。微分すると
$P' = Q + (x - \alpha)Q'$ を得る。したがって $\alpha$ は $Q$ の根である。
よって $P$ と $P'$ が共通根をもつならば、$P$ の根は相異ならない。
逆に、$P$ が重根をもつ、すなわち $(x - \alpha)^2$ が $P$ を割るならば、
$\alpha$ は $P$ と $P'$ の両方の根である。補題
\ref{fields-lemma-relatively-prime-polynomials} と合わせれば補題が証明される。
\end{proof}

\begin{lemma}
\label{fields-lemma-nr-roots-unchanged}
% P01-E0249
$F$ を体とし、$\overline{F}$ を $F$ の代数閉包とする。$F$ の標数を
$p > 0$ とし、$P$ を $F$ 上の多項式とする。このとき $\overline{F}$ における
$P$ の根の集合と $P(x^p)$ の根の集合は同じ濃度をもつ（重複度は数えない）。
\end{lemma}

\begin{proof}
明らかに、$\alpha$ が $P(x^p)$ の根であることと $\alpha^p$ が $P$ の根で
あることは同値である。言い換えれば、$P(x^p)$ の根は、$\beta$ を $P$ の根と
したときの $x^p - \beta$ の根である。したがって写像
$\overline{F} \to \overline{F}$、$\alpha \mapsto \alpha^p$ が全単射であることを
示せば十分である。$\overline{F}$ は代数閉で、すべての元が $p$ 乗根をもつので、
この写像は全射である。また標数が $p$ であるため、$\alpha^p = \beta^p$ ならば
$(\alpha - \beta)^p = 0$ となるので単射である。もちろん体では
$x^p = 0$ ならば $x = 0$ である。
\end{proof}

\noindent
$F$ を体とし、$P$ を $F$ 上の既約多項式とする。このとき、ある分離的な
既約多項式 $Q$ により $P = Q(x^q)$ であることを知っている（補題
\ref{fields-lemma-irreducible-polynomials}）。ここで $q$ は標数 $p$ の冪である
（標数が零ならば、$q = 1$\footnote{本章では $0^0 = 1$ と定めるのが
便利である。} かつ $Q = P$ とする）。補題
\ref{fields-lemma-nr-roots-unchanged} により、$F$ の任意の代数閉包における
$P$ と $Q$ の根の個数は同じである。補題 \ref{fields-lemma-recognize-separable}
により、この個数は $Q$ の次数に等しい。

\begin{definition}
\label{fields-definition-separable-degree}
$F$ を体とし、$P$ を $F$ 上の既約多項式とする。$P$ の {\it 分離次数} とは、
$F$ の任意の代数閉包における $P$ の根の集合の濃度をいう（上の議論を
参照せよ）。記号は $\deg_s(P)$ とする。
\end{definition}

\noindent
$P$ の分離次数は常にその次数を割り、その商は標数の冪である。標数が零ならば
$\deg_s(P) = \deg(P)$ である。

\begin{situation}
\label{fields-situation-finitely-generated}
% P01-E0250
ここで $F$ を体とし、$K/F$ を元 $\alpha_1, \ldots, \alpha_n \in K$ で
生成される有限次拡大とする。$K_0 = F$ と置き、
$$
K_i = F(\alpha_1, \ldots, \alpha_i)
$$
と定めると、有限次拡大の塔
$K = K_n / K_{n - 1} / \ldots / K_0 = F$
を得る。$P_i$ を $K_{i - 1}$ 上の $\alpha_i$ の最小多項式とする。
最後に、$F$ の代数閉包 $\overline{F}$ を一つ固定する。
\end{situation}

\noindent
$F$、$K$、$\alpha_i$、および $\overline{F}$ は Situation
\ref{fields-situation-finitely-generated} のものとする。
$\varphi : K \to \overline{F}$ を $F$ の拡大の射と仮定する。このとき写像
$\varphi_i : K_i \to \overline{F}$ を得る。$\varphi_{i - 1}$ による
$P_i \in K_{i - 1}[x]$ の像は多項式
$P_i^{\varphi_{i - 1}} \in \overline{F}[x]$ であり、
$\varphi(\alpha_i)$ はこの多項式の根である。

\begin{lemma}
\label{fields-lemma-count-embeddings}
Situation \ref{fields-situation-finitely-generated} において、対応
$$
\Mor_F(K, \overline{F})
\longrightarrow
\{\text{次の条件を満たす }(\beta_1, \ldots, \beta_n)\},
\quad
\varphi \longmapsto (\varphi(\alpha_1), \ldots, \varphi(\alpha_n))
$$
は全単射である。ここで右辺は、次の性質をもつ $\overline{F}$ の元の
$n$ 組 $(\beta_1, \ldots, \beta_n)$ の集合である。
\begin{enumerate}
\item $\beta_1 \in \overline{F}$ は $P_1$ の根である。
$\varphi_1 : K_1 \to \overline{F}$ を、$\alpha_1$ を $\beta_1$ に送る
$F$ 上の準同型とする。
\item $\beta_2 \in \overline{F}$ は $P_2^{\varphi_1}$ の根である。
$\varphi_2 : K_2 \to \overline{F}$ を、$\varphi_1$ を延長して $\alpha_2$ を
$\beta_2$ に送る準同型とする。
\item 以下同様に続け、最後に、
\item $\beta_n \in \overline{F}$ は $P_n^{\varphi_{n - 1}}$ の根である。
\end{enumerate}
% P01-E0251
各段階で準同型 $\varphi_i$ が存在して一意であるのは、
$K_i = K_{i - 1}[x]/(P_i)$ であり、$\beta_i$ が
$P_i^{\varphi_{i - 1}}$ の根だからである。
\end{lemma}

\begin{proof}
左辺から右辺への写像については補題の前で論じた。
$(\beta_1, \ldots, \beta_n)$ を右辺の元とする。このとき
$\varphi : K = K_n \to \overline{F}$ を、$\varphi_{n - 1}$ を延長して
$\alpha_n$ を $\beta_n$ に送る一意な準同型とする。一意性により、
二つの構成は互いに逆である。
\end{proof}

\begin{lemma}
\label{fields-lemma-count-embeddings-explicitly}
Situation \ref{fields-situation-finitely-generated} において
$|\Mor_F(K, \overline{F})| = \prod_{i = 1}^n \deg_s(P_i)$ である。
\end{lemma}

\begin{proof}
補題 \ref{fields-lemma-count-embeddings} から直ちに従う。ここで暗黙に用いている
重要な要素は、定義 \ref{fields-definition-separable-degree} の直前で確認した、
既約多項式の分離次数が良定義であるという事実である。
\end{proof}

\noindent
上の結果を用いて、分離的体拡大を特徴づける。

\begin{lemma}
\label{fields-lemma-separably-generated-separable}
仮定と記号は Situation \ref{fields-situation-finitely-generated} のものとする。
各 $P_i$ が分離的、すなわち各 $\alpha_i$ が $K_{i - 1}$ 上分離的ならば、
$$
|\Mor_F(K, \overline{F})| = [K : F]
$$
であり、体拡大 $K/F$ は分離的である。$\alpha_i$ のうち一つでも
$K_{i - 1}$ 上分離的でなければ、
$|\Mor_F(K, \overline{F})| < [K : F]$ である。
\end{lemma}

\begin{proof}
$\alpha_i$ が $K_{i - 1}$ 上分離的ならば、
$\deg_s(P_i) = \deg(P_i) = [K_i : K_{i - 1}]$
である（最後の等式は補題 \ref{fields-lemma-degree-minimal-polynomial} による）。
乗法性（補題 \ref{fields-lemma-multiplicativity-degrees}）により、
$$
[K : F] = \prod [K_i : K_{i - 1}] = \prod \deg(P_i) =
\prod \deg_s(P_i) = |\Mor_F(K, \overline{F})|
$$
となる。最後の等式は補題 \ref{fields-lemma-count-embeddings-explicitly} による。
まったく同じ議論により、$\alpha_i$ のうち一つでも $K_{i - 1}$ 上分離的で
なければ、真の不等式 $|\Mor_F(K, \overline{F})| < [K : F]$ を得る。

\medskip\noindent
再び、各 $\alpha_i$ が $K_{i - 1}$ 上分離的であると仮定する。
$K/F$ が分離的であることを示す。$\gamma = \gamma_1 \in K$ を任意に取る。
すると、$K = F(\gamma_1, \ldots, \gamma_m)$ となるような追加の元
$\gamma_2, \ldots, \gamma_m$ を見つけられる（例えば
$\gamma_2 = \alpha_1, \ldots, \gamma_{n + 1} = \alpha_n$ と取ればよい）。
% P01-E0252
このとき、補題の最後の部分（上ですでに証明済み）から、$\gamma$ が $F$ 上
分離的でなければ真の不等式
$|\Mor_F(K, \overline{F})| < [K : F]$ を得る。しかしこれは補題の最初の部分
（これも上ですでに証明済み）に反する。
\end{proof}

\begin{lemma}
\label{fields-lemma-separable-equality}
$K/F$ を体の有限次拡大とし、$\overline{F}$ を $F$ の代数閉包とする。このとき
$$
|\Mor_F(K, \overline{F})| \leq [K : F]
$$
であり、等号が成り立つことと $K$ が $F$ 上分離的であることは同値である。
\end{lemma}

\begin{proof}
これは補題 \ref{fields-lemma-separably-generated-separable} の系である。実際、
$K/F$ は有限次なので、$F$ 上 $K$ を生成する有限個の元
$\alpha_1, \ldots, \alpha_n \in K$ を見つけられる（例えば $\alpha_i$ として
$F$ 上の $K$ の基底を選べばよい）。$K/F$ が分離的ならば、補題
\ref{fields-lemma-separable-goes-up} により、各 $\alpha_i$ は
$F(\alpha_1, \ldots, \alpha_{i - 1})$ 上分離的であり、補題
\ref{fields-lemma-separably-generated-separable} により等号を得る。
一方、等号が成り立つならば、$\alpha_1, \ldots, \alpha_n$ をどのように
選んでも、補題 \ref{fields-lemma-separably-generated-separable} により
$\alpha_1$ は $F$ 上分離的である。列を $K$ の任意の元から始められるので、
$K$ は $F$ 上分離的である。
\end{proof}

\begin{lemma}
\label{fields-lemma-separable-permanence}
$E/k$ と $F/E$ を分離代数拡大とする。このとき $F/k$ は分離的体拡大である。
\end{lemma}

\begin{proof}
$\alpha \in F$ を選ぶ。このとき $\alpha$ は $E$ 上分離代数的である。
$P = x^d + \sum_{i < d} a_i x^i$ を $E$ 上の $\alpha$ の最小多項式とする。
各 $a_i$ は $k$ 上分離代数的である。体の塔
$$
k \subset k(a_0) \subset k(a_0, a_1) \subset \ldots \subset
k(a_0, \ldots, a_{d - 1}) \subset k(a_0, \ldots, a_{d - 1}, \alpha)
$$
を考える。$a_i$ は $k$ 上分離代数的なので、補題
\ref{fields-lemma-separable-goes-up} により
$k(a_0, \ldots, a_{i - 1})$ 上分離代数的である。最後に、$\alpha$ は
$k(a_0, \ldots, a_{d - 1})$ 上分離代数的である。実際、その元は $P$ の根で、
この多項式は既約（より大きい可能性のある体 $E$ 上既約だから）かつ分離的
（$E$ 上分離的だから）である。したがって補題
\ref{fields-lemma-separably-generated-separable} により
$k(a_0, \ldots, a_{d - 1}, \alpha)$ は $k$ 上分離的であり、所望のとおり
$\alpha$ が $k$ 上分離的であると結論する。
\end{proof}

\begin{lemma}
\label{fields-lemma-separable-elements}
$E/k$ を体拡大とする。このとき $E$ の $k$ 上分離的な元全体は
$E/k$ の部分拡大をなす。
\end{lemma}

\begin{proof}
$\alpha, \beta \in E$ を $k$ 上分離的とする。補題
\ref{fields-lemma-separable-goes-up} により、$\beta$ は $k(\alpha)$ 上分離的である。
補題 \ref{fields-lemma-separably-generated-separable}（$n = 2$、
$\alpha_1 = \alpha$、$\alpha_2 = \beta$ として適用）により、
$k(\alpha, \beta)$ は $k$ 上分離的である。
\end{proof}



\section{指標の線形独立性}
\label{fields-section-independence-characters}

\noindent
主張は次のとおりである。

\begin{lemma}
\label{fields-lemma-independence-characters}
$L$ を体とし、$G$ をモノイド、例えば群とする。$L$ を乗法によりモノイドと
みなしたとき、$\chi_1, \ldots, \chi_n : G \to L$ を相異なるモノイド準同型と
する。このとき $\chi_1, \ldots, \chi_n$ は $L$ 上線形独立である。すなわち、
すべてが零ではない $\lambda_1, \ldots, \lambda_n \in L$ に対し、ある
$g \in G$ について $\sum \lambda_i\chi_i(g) \not = 0$ となる。
\end{lemma}

\begin{proof}
$n = 1$ のとき、$e \in G$ が単位元ならば $\chi_1(e) = 1$ なので主張は
成り立つ。$n > 1$ に対して帰納法で証明する。
$\lambda_1, \ldots, \lambda_n \in L$ はすべてが零ではないと仮定する。
あるものについて $\lambda_i = 0$ ならば、$n$ に関する帰納法で結論を得る。
ある $g \in G$ について $\sum \lambda_i\chi_i(g) \not = 0$ を示したいので、
$-\lambda_n$ で割ることにより $\lambda_n = -1$ と仮定してよい。
問題が起こりうるのは、
$$
\chi_n(g) = \sum\nolimits_{i = 1, \ldots, n - 1} \lambda_i\chi_i(g)
$$
がすべての $g \in G$ について成り立つ場合だけである。$h \in G$ を固定する。
このときさらに
\begin{align*}
\chi_n(h)\chi_n(g) & = \chi_n(hg) \\
& = \sum\nolimits_{i = 1, \ldots, n - 1} \lambda_i\chi_i(hg) \\
& = \sum\nolimits_{i = 1, \ldots, n - 1} \lambda_i\chi_i(h) \chi_i(g)
\end{align*}
を得る。先の関係式に $\chi_n(h)$ を掛けて差を取ると、
$$
0 = \sum\nolimits_{i = 1, \ldots, n - 1}
\lambda_i (\chi_n(h) - \chi_i(h)) \chi_i(g)
$$
をすべての $g \in G$ について得る。$\lambda_i \not = 0$ なので、帰納法により
すべての $i$ について $\chi_n(h) = \chi_i(h)$ と結論する。上の $h$ の選択は
任意だったので、$i \leq n - 1$ について $\chi_i = \chi_n$ となる。
これは指標 $\chi_i$ が相異なるという仮定に反する。
\end{proof}

\begin{lemma}
\label{fields-lemma-sums-of-powers}
$L$ を体とする。$n \geq 1$ とし、$\alpha_1, \ldots, \alpha_n \in L$ を
$L$ の相異なる元とする。このとき
$\sum_{i = 1, \ldots, n} \alpha_i^e \not = 0$ となる $e \geq 0$ が存在する。
\end{lemma}

\begin{proof}
指標の線形独立性（補題 \ref{fields-lemma-independence-characters}）を、
モノイド準同型 $\mathbf{Z}_{\geq 0} \to L$、$e \mapsto \alpha_i^e$ に適用する。
\end{proof}

\begin{lemma}
\label{fields-lemma-independence-embeddings}
$K/F$ と $L/F$ を体拡大とし、
$\sigma_1, \ldots, \sigma_n : K \to L$ を相異なる $F$-拡大の射とする。
このとき $\sigma_1, \ldots, \sigma_n$ は $L$ 上線形独立である。すなわち、
すべてが零ではない $\lambda_1, \ldots, \lambda_n \in L$ に対し、ある
$\alpha \in K$ について $\sum \lambda_i\sigma_i(\alpha) \not = 0$ となる。
\end{lemma}

\begin{proof}
補題 \ref{fields-lemma-independence-characters} を、単元群への $\sigma_i$ の制限に
適用する。
\end{proof}

\begin{lemma}
\label{fields-lemma-finite-separable-tensor-alg-closed}
$K/F$ と $L/F$ を体拡大とし、$K/F$ は有限分離拡大、$L$ は代数閉とする。
このとき写像
$$
K \otimes_F L
\longrightarrow
\prod\nolimits_{\sigma \in \Hom_F(K, L)} L,\quad
\alpha \otimes \beta \mapsto (\sigma(\alpha)\beta)_\sigma
$$
は $L$-代数の同型である。
\end{lemma}

\begin{proof}
$F$ 上のベクトル空間としての $K$ の基底 $\alpha_1, \ldots, \alpha_n$ を選ぶ。
補題 \ref{fields-lemma-separable-equality}（および省略されたごく短い議論）により、
集合 $\Hom_F(K, L)$ は $n$ 個の元、例えば $\sigma_1, \ldots, \sigma_n$ をもつ。
特に、両辺は $L$ 上のベクトル空間として同じ次元 $n$ をもつ。したがって写像が
同型でなければ核をもつ。言い換えれば、すべてが零ではない
$\mu_j \in L$、$j = 1, \ldots, n$ が存在し、
$\sum \alpha_j \otimes \mu_j$ が核に属することになる。すなわち、すべての
$i$ について $\sum \sigma_i(\alpha_j)\mu_j = 0$ となる。これは成分が
$\sigma_i(\alpha_j)$ である $n \times n$ 行列が可逆でないことを意味する。
したがって、すべてが零ではない $\lambda_1, \ldots, \lambda_n \in L$ で、
すべての $j$ について $\sum \lambda_i\sigma_i(\alpha_j) = 0$ となるものを
見つけられる。いま任意の元 $\alpha \in K$ は、$\beta_j \in F$ を用いて
$\alpha = \sum \beta_j \alpha_j$ と書けるので、
$$
\sum \lambda_i\sigma_i(\alpha) =
\sum \lambda_i\sigma_i(\sum \beta_j \alpha_j) =
\sum \beta_j \sum \lambda_i\sigma_i(\alpha_j) = 0
$$
を得る。これは補題 \ref{fields-lemma-independence-embeddings} に反する。
\end{proof}
\section{純非分離拡大}
\label{fields-section-purely-inseparable}

\noindent
純非分離拡大は、前節で定義した分離拡大とは反対のものである。このような拡大は
正標数でのみ現れる。

\begin{definition}
\label{fields-definition-purely-inseparable}
$F$ を標数 $p > 0$ の体とし、$K/F$ を拡大とする。
\begin{enumerate}
\item ある $p$ の冪 $q$ に対して $\alpha^q \in F$ となるとき、
元 $\alpha \in K$ は $F$ 上 {\it 純非分離的} であるという。
\item $K$ のすべての元が $F$ 上純非分離的であることと、拡大 $K/F$ が
{\it 純非分離的} であることは同値であると定める。
\end{enumerate}
体拡大 $L/M$（$M$ の標数には条件を課さない）が与えられたとき、
$L = M$ であるか、または $M$ の標数が素数 $p$ であって $L/M$ が上で
定義した意味で純非分離的ならば、その拡大は {\it 純非分離的} であるという。
\end{definition}

\noindent
純非分離拡大は必ず代数的であることに注意する。$F$ を標数 $p > 0$ の体とする。
純非分離拡大の一例は、まだ $p$ 乗根をもたない元 $t \in F$ にその根を
添加することで得られる。実際、下の補題は $P = x^p - t$ が既約であることを
示すので、
$$
K = F[x]/(P) = F[t^{1/p}]
$$
は体である。また、$K$ の任意の元
$$
a_0 + a_1t^{1/p} + \ldots + a_{p - 1}t^{(p - 1)/p},\quad a_i \in F
$$
の $p$ 乗は
$$
(a_0 + a_1t^{1/p} + \ldots + a_{p - 1}t^{(p - 1)/p})^p =
a_0^p + a_1^p t + \ldots + a_{p - 1}^pt^{p - 1} \in F
$$
に等しいので、$K$ は $F$ 上純非分離的である。この状況は
$\mathbf{F}_p$ 上の有理関数体 $\mathbf{F}_p(t)$ で生じる。

\begin{lemma}
\label{fields-lemma-take-pth-root}
$p$ を素数とし、$F$ を標数 $p$ の体とする。$t \in F$ は $F$ に
$p$ 乗根をもたない元とする。このとき多項式 $x^p - t$ は $F$ 上既約である。
\end{lemma}

\begin{proof}
これを見るため、因数分解 $x^p - t = f g$ があると仮定する。
微分すると $f' g + f g' = 0$ を得る。$f$ と $g$ の次数は $p$ より小さいので、
$f' = 0$ でも $g' = 0$ でもないことに注意する。さらに
$\deg(f') < \deg(f)$ および $\deg(g') < \deg(g)$ である。
したがって $f$ と $g$ は共通因子をもつ。よって $x^p - t$ が可約ならば、
ある既約多項式 $f$、$c \in F^*$、および $n > 1$ により
$x^p - t = c f^n$ の形である。$p$ は素数なので、これは $n = p$ かつ
$f$ が一次式であることを含意し、したがって $x^p - t$ は $F$ に根をもつ
ことになる。これは矛盾である。
\end{proof}

\noindent
標数 $p$ において $p$ 乗根を取ることが非常に重要な操作であることを
後に見る。

\begin{lemma}
\label{fields-lemma-purely-inseparable-permanence}
$E/k$ と $F/E$ を純非分離的体拡大とする。このとき $F/k$ は
純非分離的体拡大である。
\end{lemma}

\begin{proof}
$k$ の標数を $p$ とする。$\alpha \in F$ を選ぶ。このとき、ある $p$ の冪
$q$ に対して $\alpha^q \in E$ である。すると、ある $p$ の冪 $q'$ に対して
$(\alpha^q)^{q'} \in k$ となる。したがって $\alpha^{qq'} \in k$ である。
\end{proof}

\begin{lemma}
\label{fields-lemma-purely-inseparable-elements}
$E/k$ を体拡大とする。このとき $E$ の $k$ 上純非分離的な元全体は
$E/k$ の部分拡大をなす。
\end{lemma}

\begin{proof}
$p$ を $k$ の標数とする。$\alpha, \beta \in E$ を $k$ 上純非分離的とする。
ある $p$ の冪 $q, q'$ に対して $\alpha^q \in k$ および
$\beta^{q'} \in k$ であるとする。$q''$ が $p$ の冪ならば、
$(\alpha + \beta)^{q''} = \alpha^{q''} + \beta^{q''}$ である。
したがって $q'' \geq q, q'$ ならば、$\alpha + \beta$ は $k$ 上
純非分離的であると結論する。$\alpha$ と $\beta$ の差、積、および商についても
同様である。
\end{proof}

\begin{lemma}
\label{fields-lemma-finite-purely-inseparable}
$E/F$ を標数 $p > 0$ の有限純非分離体拡大とする。このとき、
体の塔
$$
E = F(\alpha_1, \ldots, \alpha_n) \supset
F(\alpha_1, \ldots, \alpha_{n - 1}) \supset
\ldots
\supset F(\alpha_1) \supset F
$$
を得るような元の列 $\alpha_1, \ldots, \alpha_n \in E$ が存在する。
各中間拡大は次数 $p$ で、$p$ 乗根の添加によって得られる。すなわち、
$i = 1, \ldots, n$ に対して
$\alpha_i^p \in F(\alpha_1, \ldots, \alpha_{i - 1})$ は
$F(\alpha_1, \ldots, \alpha_{i - 1})$ に $p$ 乗根をもたない元である。
\end{lemma}

\begin{proof}
$E/F$ の次数に関する帰納法による。拡大の次数が $1$ ならば、結果は明らかで
ある（$n = 0$ とする）。そうでなければ、
$\alpha \in E$、$\alpha \not \in F$ を選ぶ。ある $r > 0$ に対して
$\alpha^{p^r} \in F$ とする。$r$ を最小に選び、$\alpha$ を
$\alpha^{p^{r - 1}}$ で置き換える。このとき $\alpha \not \in F$ だが
$\alpha^p \in F$ である。すると $t = \alpha^p$ は $F$ における
$p$ 乗ではない（そうであれば $\alpha \in F$ となる。補題
\ref{fields-lemma-nr-roots-unchanged} またはその証明を参照せよ）。
したがって $F \subset F(\alpha)$ は次数 $p$ の部分拡大である
（補題 \ref{fields-lemma-take-pth-root}）。帰納法により、補題の結論を満たしつつ
$E/F(\alpha)$ を生成する $\alpha_1, \ldots, \alpha_n \in E$ を見つけられる。
列 $\alpha, \alpha_1, \ldots, \alpha_n$ が拡大 $E/F$ に対する所望の列である。
\end{proof}

\begin{lemma}
\label{fields-lemma-separable-first}
\begin{slogan}
任意の代数的体拡大は、一意的に、分離的体拡大に純非分離的体拡大を続けた
ものとなる。
\end{slogan}
$E/F$ を代数的体拡大とする。$E_{sep}/F$ が分離的であり、
$E/E_{sep}$ が純非分離的となる一意な部分拡大 $E/E_{sep}/F$ が存在する。
\end{lemma}

\begin{proof}
標数が零ならば $E_{sep} = E$ と置く。標数が $p > 0$ であると仮定する。
$E_{sep}$ を、$F$ 上分離的な $E$ の元全体とする。補題
\ref{fields-lemma-separable-elements} によりこれは部分拡大であり、もちろん
$E_{sep}$ は $F$ 上分離的である。$E$ の $\alpha$ が与えられると、
$\alpha^q$ が $F$ 上分離的となる $p$ の冪 $q$ が存在する。すなわち、$q$ は、
% P01-E1267
$\alpha$ の最小多項式が $P(x^q)$ の形となり、$P$ が分離代数的となるような
$p$ の冪である。補題 \ref{fields-lemma-irreducible-polynomials} を参照せよ。
したがって $E/E_{sep}$ は純非分離的である。一意性は明らかである。
\end{proof}

\begin{definition}
\label{fields-definition-insep-degree}
$E/F$ を代数的体拡大とし、$E_{sep}$ を補題
\ref{fields-lemma-separable-first} で得られた部分拡大とする。
% P01-E1268
\begin{enumerate}
\item 整数 $[E_{sep} : F]$ をこの拡大の {\it 分離次数} という。
記号は $[E : F]_s$ とする。
\item 整数 $[E : E_{sep}]$ をこの拡大の {\it 非分離次数}、または
{\it 非分離度} という。記号は $[E : F]_i$ とする。
\end{enumerate}
\end{definition}

\noindent
もちろん標数 $0$ では $[E : F] = [E : F]_s$ および $[E : F]_i = 1$ である。
乗法性（補題 \ref{fields-lemma-multiplicativity-degrees}）により、
$$
[E : F] = [E : F]_s [E : F]_i
$$
が成り立つ。これは、これらの次数の一部が無限の場合にも成り立つ。実際、
分離次数と非分離次数も乗法的である（補題
\ref{fields-lemma-multiplicativity-all-degrees} を参照せよ）。

\begin{lemma}
\label{fields-lemma-separable-degree}
$K/F$ を有限次拡大とし、$\overline{F}$ を $F$ の代数閉包とする。このとき
$[K : F]_s = |\Mor_F(K, \overline{F})|$ である。
\end{lemma}

\begin{proof}
まず $K/F$ が純非分離的な場合を証明する。すなわち、この場合には一意な写像
$K \to \overline{F}$ があると主張する。これは補題
\ref{fields-lemma-finite-purely-inseparable} のような元の列
$\alpha_1, \ldots, \alpha_n \in K$ を選ぶことで分かる。
$F(\alpha_1, \ldots, \alpha_{i - 1})$ 上の $\alpha_i$ の既約多項式は
$x^p - \alpha_i^p$ である。補題 \ref{fields-lemma-count-embeddings-explicitly} を
適用すると $|\Mor_F(K, \overline{F})| = 1$ が分かる。一方、この場合
$[K : F]_s = 1$ なので等式が成り立つ。

\medskip\noindent
一般の有限次拡大 $K/F$ に戻る。この場合、補題
\ref{fields-lemma-separable-first} のように $F \subset K_s \subset K$ を選ぶ。
補題 \ref{fields-lemma-separable-equality} により
$|\Mor_F(K_s, \overline{F})| = [K_s : F] = [K : F]_s$ である。
一方、前段落の結果により、任意の体写像
$\sigma' : K_s \to \overline{F}$ は一意な体写像
$\sigma : K \to \overline{F}$ へ延長される。言い換えれば、
$|\Mor_F(K, \overline{F})| = |\Mor_F(K_s, \overline{F})|$ であり、
証明が完了する。
\end{proof}

\begin{lemma}[乗法性]
\label{fields-lemma-multiplicativity-all-degrees}
代数的体拡大の塔 $K/E/F$ が与えられたと仮定する。このとき
$$
[K : F]_s = [K : E]_s [E : F]_s
\quad\text{かつ}\quad
[K : F]_i = [K : E]_i [E : F]_i
$$
である。
\end{lemma}

\begin{proof}
まず $K$ が $F$ 上有限次の場合に証明する。通常の次数については乗法性
（補題 \ref{fields-lemma-multiplicativity-degrees}）があるので、二つの式の一方を
証明すれば十分である。補題 \ref{fields-lemma-separable-degree} により
$[K : F]_s = |\Mor_F(K, \overline{F})|$ である。同じ補題により、任意の
$\sigma \in \Mor_F(E, \overline{F})$ が与えられたとき、$\sigma$ を写像
$\tau : K \to \overline{F}$ へ延長する仕方は $[K : E]_s$ 個ある。実際、
$E \cong \sigma(E) \subset \overline{F}$ を通じて $\overline{F}$ を
$E$ の代数閉包とみなせる。$\sigma$ の選択肢が
$[E : F]_s = |\Mor_F(E, \overline{F})|$ 個あるという事実と合わせれば
結果を得る。

\medskip\noindent
拡大が無限の場合の証明は省略する。
\end{proof}
\section{正規拡大}
\label{fields-section-normal}

\noindent
$P \in F[x]$ を体 $F$ 上の非定数多項式とする。ある
$c \in F^*$、$n \geq 1$、$\alpha_1, \ldots, \alpha_n \in F$ により
$$
P = c(x - \alpha_1) \ldots (x - \alpha_n)
$$
が $F[x]$ で成り立つとき、$P$ は {\it $F$ 上一次因子に完全分解する}、
または {\it $F$ 上完全分解する} という。正規拡大を次のように定義する。

\begin{definition}
\label{fields-definition-normal}
$E/F$ を代数的体拡大とする。任意の $\alpha \in E$ に対し、$F$ 上の
$\alpha$ の最小多項式 $P$ が $E$ 上一次因子に完全分解するとき、
$E$ は $F$ 上 {\it 正規} であるという。
\end{definition}

\noindent
分離拡大の場合と同様、この概念の基本的性質を確立するには多少の議論が必要で
ある。

\begin{lemma}
\label{fields-lemma-normal-goes-up}
$K/E/F$ を代数的体拡大の塔とする。$K$ が $F$ 上正規ならば、$K$ は
$E$ 上正規である。
\end{lemma}

\begin{proof}
$\alpha \in K$ とする。$P$ を $F$ 上の $\alpha$ の最小多項式とし、$Q$ を
$E$ 上の $\alpha$ の最小多項式とする。このとき $Q$ は多項式環 $E[x]$ で
$P$ を割り、$P = QR$ と書ける。したがって $P$ が $K$ 上完全分解するならば、
$Q$ もそうである。
\end{proof}

\begin{lemma}
\label{fields-lemma-intersect-normal}
$F$ を体とし、$M/F$ を代数拡大とする。$M/E_i/F$、$i \in I$ を
$E_i/F$ が正規である部分拡大とする。このとき $\bigcap E_i$ は $F$ 上
正規である。
\end{lemma}

\begin{proof}
定義から直ちに従う。
\end{proof}

\begin{lemma}
\label{fields-lemma-separable-first-normal}
$E/F$ を正規代数的体拡大とする。このとき、補題
\ref{fields-lemma-separable-first} の部分拡大 $E/E_{sep}/F$ は正規である。
\end{lemma}

\begin{proof}
標数が零ならば $E_{sep} = E$ であり、結果は明らかである。標数が
$p > 0$ ならば、$E_{sep}$ は $F$ 上分離的な $E$ の元全体である。
このとき $\alpha \in E_{sep}$ の最小多項式を $P$ とし、
$P = c(x - \alpha)(x - \alpha_2) \ldots (x - \alpha_d)$ と書く。
ここで $\alpha_2, \ldots, \alpha_d \in E$ である。$P$ は分離多項式であり、
$\alpha_i$ は $P$ の根なので、所望のとおり $\alpha_i \in E_{sep}$ と
結論する。
\end{proof}

\begin{lemma}
\label{fields-lemma-characterize-normal}
$E/F$ を体の代数拡大とし、$\overline{F}$ を $F$ の代数閉包とする。
次の条件は同値である。
\begin{enumerate}
\item $E$ は $F$ 上正規である。
\item 任意の組 $\sigma, \sigma' \in \Mor_F(E, \overline{F})$ に対して
$\sigma(E) = \sigma'(E)$ である。
\end{enumerate}
\end{lemma}

\begin{proof}
$\mathcal{P}$ を、$E$ のすべての元の $F$ 上の最小多項式全体の集合とする。
$$
T =
\{\beta \in \overline{F} \mid P(\beta) = 0\text{ を満たす }P \in \mathcal{P}\text{ が存在する}\}
$$
と置く。$E$ が $F$ 上正規ならば、任意の
$\sigma \in \Mor_F(E, \overline{F})$ に対して $\sigma(E) = T$ であることは
明らかである。したがって (1) は (2) を含意する。

\medskip\noindent
逆に (2) を仮定し、$\beta \in T$ を選ぶ。対応する $\alpha \in E$ で、その
最小多項式 $P \in \mathcal{P}$ が $\beta$ を零にするものを見つけられる。
$F(\alpha) = F[x]/(P)$ なので、$\alpha$ を $\beta$ に送る元
$\sigma_0 \in \Mor_F(F(\alpha), \overline{F})$ を見つけられる。
補題 \ref{fields-lemma-map-into-algebraic-closure} により、$\sigma_0$ を
$\sigma \in \Mor_F(E, \overline{F})$ へ延長できる。よって $\beta$ はすべての
埋め込み $\sigma : E \to \overline{F}$ の共通の像に属する。したがって任意の
$\sigma$ に対し $\sigma(E) = T$ である。一つの $\sigma$ を固定する。
いま $P \in \mathcal{P}$ とする。補題 \ref{fields-lemma-algebraically-closed} により、
ある $n$ と $\beta_i \in \overline{F}$ に対して
$$
P = (x - \beta_1) \ldots (x - \beta_n)
$$
と書ける。$\beta_i \in T$ であることに注意する。したがって、ある
$\alpha_i \in E$ により $\beta_i = \sigma(\alpha_i)$ である。よって
$P = (x - \alpha_1) \ldots (x - \alpha_n)$ は $E$ 上完全分解する。
これで証明が完了する。
\end{proof}

\begin{lemma}
\label{fields-lemma-normally-generated}
$E/F$ を体の代数拡大とする。$E$ が $F$ 上の $\alpha_i \in E$、$i \in I$ で
生成され、各 $i$ について $F$ 上の $\alpha_i$ の最小多項式が $E$ で
完全分解するならば、$E/F$ は正規である。
\end{lemma}

\begin{proof}
$P_i$ を $F$ 上の $\alpha_i$ の最小多項式とする。
$\alpha_i = \alpha_{i, 1}, \alpha_{i, 2}, \ldots, \alpha_{i, d_i}$ を
$E$ 上の $P_i$ の根とする。$F$ 上の二つの埋め込み
$\sigma, \sigma' : E \to \overline{F}$ が与えられたとき、
$$
\{\sigma(\alpha_{i, 1}), \ldots, \sigma(\alpha_{i, d_i})\} =
\{\sigma'(\alpha_{i, 1}), \ldots, \sigma'(\alpha_{i, d_i})\}
$$
である。実際、両辺は $\overline{F}$ における $P_i$ の根の集合に等しい。
元 $\alpha_{i, j}$ は $F$ 上 $E$ を生成するので、
$\sigma(E) = \sigma'(E)$ である。したがって補題
\ref{fields-lemma-characterize-normal} により $E/F$ は正規である。
\end{proof}

\begin{lemma}
\label{fields-lemma-lift-maps}
$L/M/K$ を代数拡大の塔とする。
\begin{enumerate}
\item $M/K$ が正規ならば、$L/K$ の任意の自己同型 $\tau$ は自己同型
$\tau|_M : M \to M$ を誘導する。
\item $L/K$ が正規ならば、任意の $K$-代数写像 $\sigma : M \to L$ は
$L$ の自己同型へ延長される。
\end{enumerate}
\end{lemma}

\begin{proof}
$L$ の代数閉包 $\overline{L}$ を選ぶ
（定理 \ref{fields-theorem-existence-algebraic-closure}）。

\medskip\noindent
$\tau$ は (1) のものとする。補題 \ref{fields-lemma-characterize-normal} により、
$\overline{L}$ の部分体として $\tau(M) = M$ である。したがって
$\tau|_M : M \to M$ は自己同型である。

\medskip\noindent
$\sigma : M \to L$ は (2) のものとする。補題
\ref{fields-lemma-map-into-algebraic-closure} により、$\sigma$ を写像
$\tau : L \to \overline{L}$ へ延長できる。すなわち、
$$
\xymatrix{
L \ar[r]_\tau & \overline{L} \\
M \ar[u] \ar[ru]_\sigma & K \ar[l] \ar[u]
}
$$
は可換である。補題 \ref{fields-lemma-characterize-normal} により
$\tau(L) = L$ である。したがって $\tau : L \to L$ は $\sigma$ を延長する
自己同型である。
\end{proof}

\begin{definition}
\label{fields-definition-automorphisms}
$E/F$ を体拡大とする。$\text{Aut}(E/F)$ または $\text{Aut}_F(E)$ は、
$F$-拡大の圏の対象としての $E$ の自己同型群を表す。
$\text{Aut}(E/F)$ の元を {\it $F$ 上の $E$ の自己同型}、または
{\it $E/F$ の自己同型} という。
\end{definition}

\noindent
自己同型による正規拡大の特徴づけを与える。

\begin{lemma}
\label{fields-lemma-normal-and-automorphisms}
$E/F$ を有限次拡大とする。このとき
$$
|\text{Aut}(E/F)| \leq [E : F]_s
$$
であり、等号が成り立つことと $E$ が $F$ 上正規であることは同値である。
\end{lemma}

\begin{proof}
$F$ の代数閉包 $\overline{F}$ を選ぶ。
$[E : F]_s = |\Mor_F(E, \overline{F})|$ であったことを思い出す。
元 $\sigma_0 \in \Mor_F(E, \overline{F})$ を選ぶ。このとき写像
$$
\text{Aut}(E/F) \longrightarrow \Mor_F(E, \overline{F}),\quad
\tau \longmapsto \sigma_0 \circ \tau
$$
は単射である。したがって不等式が成り立つ。等号が成り立つならば、任意の
$\sigma \in \Mor_F(E, \overline{F})$ は $\sigma_0$ に自己同型を前合成することで
得られる。よって $\sigma(E) = \sigma_0(E)$ である。補題
\ref{fields-lemma-characterize-normal} により $E$ は $F$ 上正規である。

\medskip\noindent
逆に $E/F$ が正規であると仮定する。すると補題
\ref{fields-lemma-characterize-normal} により、任意の
$\sigma \in \Mor_F(E, \overline{F})$ に対して
$\sigma(E) = \sigma_0(E)$ である。したがって
$\tau = \sigma_0^{-1} \circ \sigma$ と置けば、$F$ 上の $E$ の自己同型を得る。
よって上に表示した写像は全射である。
\end{proof}

\begin{lemma}
\label{fields-lemma-normal-embeddings-differ-by-aut}
$L/K$ を体の代数的正規拡大とし、$E/K$ を体拡大とする。このとき、
$L$ から $E$ への $K$-埋め込みが存在しないか、または一つの
$\tau : L \to E$ が存在し、他のすべては
$\sigma \in \text{Aut}(L/K)$ を用いて $\tau \circ \sigma$ の形である。
\end{lemma}

\begin{proof}
$\tau$ が与えられたとき、$L$ を $\tau(L) \subset E$ で置き換え、
補題 \ref{fields-lemma-lift-maps} を適用する。
\end{proof}
\section{分解体}
% P01-E0253
\label{fields-section-splitting-fieds}

\noindent
次の補題は正規体拡大を構成するための有用な道具である。

\begin{lemma}
\label{fields-lemma-splitting-field}
$F$ を体とし、$P \in F[x]$ を非定数多項式とする。$P$ が $E$ 上完全分解する
ような最小の体拡大 $E/F$ が存在する。さらに体拡大 $E/F$ は正規であり、
（一意とは限らない）同型を除いて一意である。
\end{lemma}

\begin{proof}
代数閉包 $\overline{F}$ を選ぶ。補題 \ref{fields-lemma-algebraically-closed} により、
$\overline{F}[x]$ において
$P = c (x - \beta_1) \ldots (x - \beta_n)$ と書ける。
$c \in F^*$ であることに注意する。
$E = F(\beta_1, \ldots, \beta_n)$ と置く。このとき、$P$ が $E$ 上完全分解する
という要件のもとで $E$ が最小であることは明らかである。

\medskip\noindent
次に、$P$ が $E'$ 上完全分解するような $F$ の別の最小体拡大を $E'$ とする。
$c \in F$ および $\alpha_i \in E'$ により
$P = c (x - \alpha_1) \ldots (x - \alpha_n)$ と書く。再び最小性から
$E' = F(\alpha_1, \ldots, \alpha_n)$ が従う。さらに、任意の
$\sigma : E' \to \overline{F}$ を選べば
（補題 \ref{fields-lemma-map-into-algebraic-closure}）、ある置換
$\tau : \{1, \ldots, n\} \to \{1, \ldots, n\}$ に対して
$\sigma(\alpha_i) = \beta_{\tau(i)}$ であることが直ちに分かる。
したがって $\sigma(E') = E$ である。補題 \ref{fields-lemma-characterize-normal}
により、これは $E'$ が $F$ の正規拡大であることを含意する。また
$F$ の拡大として $E \cong E'$ であり、証明が完了する。
\end{proof}

\begin{definition}
\label{fields-definition-splitting-field}
$F$ を体とし、$P \in F[x]$ を非定数多項式とする。補題
\ref{fields-lemma-splitting-field} で構成した体拡大 $E/F$ を
{\it $F$ 上の $P$ の分解体} という。
\end{definition}

\begin{lemma}
\label{fields-lemma-normal-closure}
\begin{slogan}
体の有限次拡大の正規閉包は存在する。
\end{slogan}
$E/F$ を体の有限次拡大とする。$K$ が $F$ 上正規となるような最小の
有限次拡大 $K/E$ が一意に存在する。
\end{lemma}

\begin{proof}
$F$ 上の $E$ の生成元 $\alpha_1, \ldots, \alpha_n$ を選ぶ。
$P_1, \ldots, P_n$ を $F$ 上の $\alpha_1, \ldots, \alpha_n$ の
最小多項式とする。$P = P_1 \ldots P_n$ と置く。各
$(x - \alpha_i)$ は $P_i$ を割るので、
$(x - \alpha_1) \ldots (x - \alpha_n)$ は $P$ を割ることに注意する。
$P = (x - \alpha_1) \ldots (x - \alpha_n)Q$ と書く。
$K/E$ を $E$ 上の $P$ の分解体とする。$K$ は $F$ 上でも $P$ の分解体である
と主張する（これにより $K$ は $F$ 上正規となる）。
$K/E$ は $E$ 上の $Q$ の根で生成され、$E$ は $F$ 上の
$(x - \alpha_1) \ldots (x - \alpha_n)$ の根で生成されるので、これは明らかで
ある。したがって $K$ は $F$ 上の $P$ の根で生成される。

\medskip\noindent
一意性を示す。$K'/E$ を、$K'/F$ が正規となる別の最小拡大とする。
代数閉包 $\overline{F}$ と埋め込み $\sigma_0 : E \to \overline{F}$ を選ぶ。
補題 \ref{fields-lemma-map-into-algebraic-closure} により、$\sigma_0$ を
$\sigma : K \to \overline{F}$ および
$\sigma' : K' \to \overline{F}$ へ延長できる。補題
\ref{fields-lemma-intersect-normal} により
$\sigma(K) \cap \sigma'(K')$ は $F$ 上正規である。最小性から
$\sigma(K) = \sigma'(K')$ と結論する。したがって
$\sigma^{-1} \circ \sigma' : K' \to K$ は $E$ の拡大の同型を与える。
\end{proof}

\begin{definition}
\label{fields-definition-normal-closure}
$E/F$ を体の有限次拡大とする。補題 \ref{fields-lemma-normal-closure} で構成した
体拡大 $K/E$ を
% P01-E0254
{\it $F$ 上の $E$ の正規閉包} という。
\end{definition}

\noindent
与えられた任意の正規拡大の内部で正規閉包を構成できる。

\begin{lemma}
\label{fields-lemma-normal-closure-inside-normal}
$L/K$ を代数的正規拡大とする。
\begin{enumerate}
\item $L/M/K$ が $M/K$ を有限次とする部分拡大ならば、
$M'/K$ が有限次かつ正規となる塔 $L/M'/M/K$ が存在する。
\item $L/M'/M/K$ が $M/K$ を正規、$M'/M$ を有限次とする塔ならば、
$M''/M$ が有限次かつ $M''/K$ が正規となる塔 $L/M''/M'/M/K$ が存在する。
\end{enumerate}
\end{lemma}

\begin{proof}
(1) を証明する。$M'$ を、$M$ を含み $K$ 上正規である $L/K$ の最小部分拡大と
する。補題 \ref{fields-lemma-normal-closure} により、これは $M/K$ の正規閉包であり、
$K$ 上有限次である。

\medskip\noindent
(2) を証明する。$\alpha_1, \ldots, \alpha_n \in M'$ を $M$ 上の $M'$ の
生成元とする。$P_1, \ldots, P_n$ を $K$ 上の
$\alpha_1, \ldots, \alpha_n$ の最小多項式とする。
$\alpha_{i, j}$ を $L$ における $P_i$ の根とする。
$M''  = M(\alpha_{i, j})$ と置く。補題 \ref{fields-lemma-normally-generated}
（生成元の集合を $M \cup \{\alpha_{i, j}\}$ として適用）により、
$M''$ は $K$ 上正規である。
\end{proof}

\noindent
次の補題は正規閉包の性質を証明するために使えることがある。

\begin{lemma}
\label{fields-lemma-normal-closure-tensor-product}
$L/K$ を有限次拡大とし、$M/L$ を $K$ 上の $L$ の正規閉包とする。このとき
全射
$$
L \otimes_K L \otimes_K \ldots \otimes_K L \longrightarrow M
$$
が $K$-代数の間に存在し、テンソル因子の個数は
$[L : K]_s \leq [L : K]$ と取れる。
\end{lemma}

\begin{proof}
$K$ の代数閉包 $\overline{K}$ を選ぶ。補題 \ref{fields-lemma-separable-degree} による
等式を用いて $n = [L : K]_s = |\Mor_K(L, \overline{K})|$ と置く。
$\Mor_K(L, \overline{K}) = \{\sigma_1, \ldots, \sigma_n\}$ とする。
$M' \subset \overline{K}$ を、$\sigma_i(L)$、$i = 1, \ldots, n$ で生成される
$K$-部分代数とする。$\overline{K}$ の任意の $K$-部分代数は体なので、$M'$ は
体である。$M'$ から $\overline{K}$ への任意の $K$-代数写像は $\sigma_i$ を
置換し、したがって $M'$ を $M'$ の中かつ上へ写す。構成により、体 $M'$ は
$\sigma_1(L)$ の元の共役で生成される。以上から補題
\ref{fields-lemma-characterize-normal} により、$M'$ は $K$ 上正規であり、
$\sigma_1(L)$ を含む $\overline{K}$ の最小正規部分拡大である。
正規閉包の一意性により $M \cong M'$ である。最後に、全射
$$
L \otimes_K L \otimes_K \ldots \otimes_K L \longrightarrow M',
\quad
\lambda_1 \otimes \ldots \otimes \lambda_n \longmapsto
\sigma_1(\lambda_1) \ldots \sigma_n(\lambda_n)
$$
が存在する。定義により $n \leq [L : K]$ であることにも注意する。
\end{proof}
\section{1 の冪根}
\label{fields-section-roots-of-1}

\noindent
$F$ を体とする。整数 $n \geq 1$ に対して
$$
\mu_n(F) = \{\zeta \in F \mid \zeta^n = 1\}
$$
と置く。これを {\it 1 の $n$ 乗根の群}、または
{\it $1$ の $n$ 乗根} という。これは乗法に関する可換群で、単位元は
$1$ である。体では次数 $d$ の多項式の根の個数は常に高々 $d$ であることに
注意する。したがって、この集合は次数 $n$ の多項式方程式で定義されるので、
$|\mu_n(F)| \leq n$ である。もちろん $\mu_n(F)$ の各元の位数は $n$ を割る。
さらに、部分群
$$
\mu_d(F) \subset \mu_n(F),\quad d | n
$$
はそれぞれ高々 $d$ 個の元をもつ。これにより $\mu_n(F)$ は巡回群である。

\begin{lemma}
\label{fields-lemma-cyclic}
$A$ を指数が $n$ を割る可換群とし、すべての $d | n$ に対して
$\{x \in A \mid dx = 0\}$ の濃度が高々 $d$ であるとする。
このとき $A$ は位数が $n$ を割る巡回群である。
\end{lemma}

\begin{proof}
条件から $|A| \leq n$ となり、特に $A$ は有限である。有限可換群の構造により、
ある整数 $1 < e_1 | e_2 | \ldots | e_r$ に対して
$A = \mathbf{Z}/e_1\mathbf{Z} \oplus \ldots \oplus \mathbf{Z}/e_r\mathbf{Z}$
である。これは $\{x \in A \mid e_1 x = 0\}$ の濃度が $e_1^r$ であることを
含意する。したがって $r = 1$ である。
\end{proof}

\noindent
これを体 $\mathbf{F}_p$ に適用すると、群
$(\mathbf{Z}/p\mathbf{Z})^*$ が巡回群であるという有名な結果を得る。
これについては有限体の節でさらに述べる。

\medskip\noindent
もう一つ、しばしば有用な観察がある。$F$ の標数が $p > 0$ ならば、
$\mu_{p^n}(F) = \{1\}$ である。実際、補題
\ref{fields-lemma-nr-roots-unchanged} の証明で見たように、標数 $p$ の体では
$p$ 乗する写像は単射だからである。（もちろん、これは補題そのものの主張からも
従う。）





\section{有限体}
\label{fields-section-finite}

\noindent
$F$ を有限体とする。単射 $\mathbf{Q} \to F$ は存在しえないので、$F$ の標数が
正であることは明らかである。$F$ の標数を $p$ とする。拡大
$\mathbf{F}_p \subset F$ は有限次である。したがって、ある $f \geq 1$ に対して
$F$ は $q = p^f$ 個の元をもつ。

\medskip\noindent
単元群 $F^*$ を考えよう。これは有限可換群なので、ある指数 $e$ をもつ。
すると $F^* = \mu_e(F)$ であり、Section \ref{fields-section-roots-of-1} の議論から
$F^*$ は位数 $q - 1$ の巡回群である。（事後的には $e = q - 1$ も従う。）
特に、$\alpha \in F^*$ が生成元ならば、明らかに
$$
F = \mathbf{F}_p(\alpha)
$$
である。言い換えれば、拡大 $F/\mathbf{F}_p$ は単一の元で生成される。
もちろん、有限体の任意の拡大 $E/F$ についても同じことが成り立つ
（$E$ はすでに素体上単一の元で生成されるからである）。





\section{原始元}
\label{fields-section-primitive-element}

\noindent
$E/F$ を体の有限次拡大とする。$E = F(\alpha)$ である元 $\alpha \in E$ を
{\it $F$ 上の $E$ の原始元} という。

\begin{lemma}[原始元]
\label{fields-lemma-primitive-element}
$E/F$ を体の有限次拡大とする。次の条件は同値である。
\begin{enumerate}
\item $F$ 上の $E$ の原始元が存在する。
\item 部分拡大 $E/K/F$ は有限個しかない。
\end{enumerate}
さらに、$E/F$ が分離的ならば (1) と (2) は成り立つ。
\end{lemma}

\begin{proof}
$\alpha \in E$ を原始元とし、$P$ を $F$ 上の $\alpha$ の最小多項式とする。
次に $E/K/F$ を部分拡大とし、$Q$ を $K$ 上の $\alpha$ の最小多項式とする。
$\deg(Q) = [E : K]$ であることに注意する。
$Q = x^d + \sum_{i < d} a_i x^i$ と書き、
$K$ が $L = F(a_0, \ldots, a_{d - 1})$ に等しいと主張する。実際、$\alpha$ の
$L$ 上の次数は $d$ であり、$L \subset K$ である。したがって
$[E : L] = [E : K]$ となり、$[K : L] = 1$、すなわち $K = L$ が従う。
よって多項式 $Q$ の可能性が有限個しかないことを示せば十分である。これは
明らかである。実際、$K[x]$、したがって特に $E[x]$ で因数分解
$P = QR$ がある。$E[x]$ では一意分解が成り立つので、$E[x]$ における
$P$ のモニック因子は有限個しかない。

\medskip\noindent
$F$ が有限体（同値なことに $E$ が有限体）ならば、Section
\ref{fields-section-finite} の議論により $E/F$ は原始元をもつ。次に $F$ が無限で、
真の部分体 $E/K/F$ が高々有限個しかないと仮定する。それらを
$K_1, \ldots, K_N$ と列挙する。各 $K_i \subset E$ は真の
$F$-ベクトル部分空間である。$F$ は無限なので、すべての $i$ に対して
$\alpha \not \in K_i$ となるベクトル $\alpha \in E$ を見つけられる
（ベクトル空間が有限個の真のベクトル部分空間の和集合に等しくなることはない。
詳細は省略する）。すると $\alpha$ は $F$ 上の $E$ の原始元である。

\medskip\noindent
(1) と (2) の同値性を確立したので、補題の最後の主張に移る。
$F$ の代数閉包 $\overline{F}$ を選び、その元
$\sigma_1, \ldots, \sigma_n \in \Mor_F(E, \overline{F})$ を列挙する。
$E/F$ は分離的なので、補題
\ref{fields-lemma-separable-equality} により $n = [E : F]$ である。
$i \not = j$ ならば、
$$
V_{ij} = \Ker(\sigma_i - \sigma_j : E \longrightarrow \overline{F})
$$
は $E$ に等しくないことに注意する。したがって前段落と同様に論じると、
すべての $i \not = j$ に対して $\alpha \not \in V_{ij}$ となる
$\alpha \in E$ を見つけられる。すると
$|\Mor_F(F(\alpha), \overline{F})| \geq n$ である。一方、
$[F(\alpha) : F] \leq [E : F]$ である。したがって補題
\ref{fields-lemma-separable-equality} により等号が成り立ち、
$E = F(\alpha)$ と結論する。
\end{proof}






\section{トレースとノルム}
\label{fields-section-trace-pairing}

\noindent
$L/K$ を体の有限次拡大とする。補題
\ref{fields-lemma-vector-space-is-free} により、$K$-加群の同型
$L \cong K^{\oplus n}$ を選べる。もちろん $n = [L : K]$ は体拡大の次数である。
この同型を用いると、$K$-代数写像
$$
L \longrightarrow \text{Mat}(n \times n, K),\quad
\alpha \longmapsto \alpha\text{ 倍写像の表現行列}
$$
を得る。したがって $\alpha \in L$ が与えられると、対応する行列のトレースと行列式を
とることができる。もちろん、これらの量は上で選んだ基底に依存しない。より標準的には、
$L$ を有限次元 $K$-ベクトル空間とみなすだけで、
$\text{Trace}_K(\alpha : L \to L)$ と行列式
$\det_K(\alpha : L \to L)$ が得られる。

\begin{definition}
\label{fields-definition-trace-norm}
$L/K$ を体の有限次拡大とする。$\alpha \in L$ に対し、{\it トレース}を
$\text{Trace}_{L/K}(\alpha) = \text{Trace}_K(\alpha : L \to L)$、
{\it ノルム}を
$\text{Norm}_{L/K}(\alpha) = \det_K(\alpha : L \to L)$
と定義する。
\end{definition}

\noindent
定義から、$\text{Trace}_{L/K}$ は $K$-線形であり、$\alpha \in K$ に対して
$\text{Trace}_{L/K}(\alpha) = [L : K]\alpha$ を満たすことは明らかである。
同様に $\text{Norm}_{L/K}$ は乗法的であり、$\alpha \in K$ に対して
$\text{Norm}_{L/K}(\alpha) = \alpha^{[L : K]}$ である。これは演習
\ref{exercises-exercise-trace-det} および
% P01-E1269
\ref{exercises-exercise-trace-det-rings} で論じる、より一般的な構成の特別な場合である。

\begin{lemma}
\label{fields-lemma-characteristic-vs-minimal-polynomial}
$L/K$ を体の有限次拡大とする。$\alpha \in L$ とし、$P$ を $K$ 上の
$\alpha$ の最小多項式とする。このとき、$K$-線形写像 $\alpha : L \to L$ の
特性多項式は $P^e$ に等しく、$e \deg(P) = [L : K]$ である。
\end{lemma}

\begin{proof}
$K(\alpha)$ 上の $L$ の基底 $\beta_1, \ldots, \beta_e$ を選ぶ。補題
\ref{fields-lemma-degree-minimal-polynomial} および
\ref{fields-lemma-multiplicativity-degrees} により、$e$ は
$e \deg(P) = [L : K]$ を満たす。すると
$L = \bigoplus K(\alpha) \beta_i$ は $\alpha$-不変部分空間への直和分解である。
したがって、$\alpha : L \to L$ の特性多項式は
$\alpha : K(\alpha) \to K(\alpha)$ の特性多項式の $e$ 乗に等しい。

\medskip\noindent
証明を完了するには $L = K(\alpha)$ と仮定してよい。この場合、
ケイリー--ハミルトンの定理により $\alpha$ は特性多項式の根である。
また特性多項式と最小多項式は同じ次数をもつので、両者は等しい。
\end{proof}

\begin{lemma}
\label{fields-lemma-trace-and-norm-from-minimal-polynomial}
$L/K$ を体の有限次拡大とする。$\alpha \in L$ とし、
$P = x^d + a_1 x^{d - 1} + \ldots + a_d$ を $K$ 上の
$\alpha$ の最小多項式とする。このとき
$$
\text{Norm}_{L/K}(\alpha) = (-1)^{[L : K]} a_d^e
\quad\text{かつ}\quad
\text{Trace}_{L/K}(\alpha) = - e a_1
$$
である。ただし $e d = [L : K]$ とする。
\end{lemma}

\begin{proof}
補題 \ref{fields-lemma-characteristic-vs-minimal-polynomial} と定義から直ちに従う。
\end{proof}

\begin{lemma}
\label{fields-lemma-trace-and-norm-linear}
$L/K$ を体の有限次拡大とする。$V$ を $L$ 上の有限次元ベクトル空間とし、
$\varphi : V \to V$ を $L$-線形写像とする。このとき
$$
\text{Trace}_K(\varphi : V \to V) =
\text{Trace}_{L/K}(\text{Trace}_L(\varphi : V \to V))
$$
および
$$
\det\nolimits_K(\varphi : V \to V) =
\text{Norm}_{L/K}(\det\nolimits_L(\varphi : V \to V))
$$
が成り立つ。
\end{lemma}

\begin{proof}
$V = L^{\oplus n}$ となる同型を選び、$\varphi$ を $n \times n$ 行列に対応させる。
トレースの場合、公式の両辺は $\varphi$ について加法的である。したがって、
$\varphi$ は $(i, j)$ 成分にちょうど一つだけ非零成分をもつ行列に対応すると
仮定してよい。この場合、直接計算すれば両辺が等しいことが分かる。

\medskip\noindent
ノルムの場合、$\varphi$ が非零な核をもてば両辺はともに零である。
したがって $\varphi$ は $\text{GL}_n(L)$ の元に対応すると仮定してよい。
公式の両辺は $\varphi$ について乗法的である。$\text{GL}_n(L)$ の各元は
基本行列の積なので、$\varphi$ は
$$
E_{12}(\lambda) =
\left(
\begin{matrix}
1 & \lambda & \ldots \\
0 & 1 & \ldots \\
\ldots & \ldots & \ldots
\end{matrix}
\right)
\quad\text{or}\quad
E_1(a) =
\left(
\begin{matrix}
a & 0 & \ldots \\
0 & 1 & \ldots \\
\ldots & \ldots & \ldots
\end{matrix}
\right)
$$
のいずれかの形であると仮定してよい（必要なら基底元を並べ替えてもよい）。
どちらの場合も、公式は直接計算で容易に確かめられる。
\end{proof}

\begin{lemma}
\label{fields-lemma-trace-and-norm-tower}
$M/L/K$ を体の有限次拡大の塔とする。このとき
$$
\text{Trace}_{M/K} = \text{Trace}_{L/K} \circ \text{Trace}_{M/L}
\quad\text{and}\quad
\text{Norm}_{M/K} = \text{Norm}_{L/K} \circ \text{Norm}_{M/L}
$$
である。
\end{lemma}

\begin{proof}
$M$ を $L$ 上のベクトル空間とみなし、補題
\ref{fields-lemma-trace-and-norm-linear} を適用する。
\end{proof}

\noindent
トレース形式はトレースを用いて定義される。

\begin{definition}
\label{fields-definition-trace-pairing}
$L/K$ を体の有限次拡大とする。$L/K$ の {\it トレース形式}とは、対称な
$K$-双線形形式
$$
Q_{L/K} : L \times L \longrightarrow K,\quad
(\alpha, \beta) \longmapsto \text{Trace}_{L/K}(\alpha\beta)
$$
のことである。
\end{definition}

\noindent
体の有限次拡大が分離的であることと、そのトレース形式が非退化であることは同値である。

\begin{lemma}
\label{fields-lemma-separable-trace-pairing}
$L/K$ を体の有限次拡大とする。次の条件は同値である。
\begin{enumerate}
\item $L/K$ は分離的である。
\item $\text{Trace}_{L/K}$ は恒等的に零ではない。
\item トレース形式 $Q_{L/K}$ は非退化である。
\end{enumerate}
\end{lemma}

\begin{proof}
(3) が (2) を含意することは明らかである。(2) が成り立つとし、
$\text{Trace}_{L/K}(\gamma) \not = 0$ となる $\gamma \in L$ を選ぶ。
すると非零な $\alpha \in L$ に対して
$Q_{L/K}(\alpha, \gamma/\alpha) \not = 0$ である。したがって
$Q_{L/K}$ は非退化である。これで (2) と (3) の同値性が示された。

\medskip\noindent
$K$ の標数が $p$ で、$\alpha \not \in K$、$\alpha^p \in K$ を満たす
$L = K(\alpha)$ であるとする。このとき
$\text{Trace}_{L/K}(1) = p = 0$ である。$i = 1, \ldots, p - 1$ に対し、
$x^p - \alpha^{pi}$ は $K$ 上の $\alpha^i$ の最小多項式であるから、補題
\ref{fields-lemma-trace-and-norm-from-minimal-polynomial} により
$\text{Trace}_{L/K}(\alpha^i) = 0$ を得る。したがって、この種の純非分離的な
次数 $p$ の拡大では、$\text{Trace}_{L/K}$ は恒等的に零である。

\medskip\noindent
$L/K$ が分離的でないと仮定する。このとき、前段落のように $L/K'$ が
純非分離的な次数 $p$ の拡大となる部分体 $L/K'/K$ が存在する。補題
\ref{fields-lemma-separable-first} および
\ref{fields-lemma-finite-purely-inseparable} を参照せよ。したがって補題
\ref{fields-lemma-trace-and-norm-tower} により、$\text{Trace}_{L/K}$ は恒等的に零である。

\medskip\noindent
一方、$L/K$ が分離的であると仮定する。次数に関する帰納法により、
$\text{Trace}_{L/K}$ が恒等的に零でないことを示す。したがって補題
\ref{fields-lemma-trace-and-norm-tower} により、$L/K$ は単一の元 $\alpha$ で生成されると
仮定してよい（トレースが非零ならば全射であることを用いる）。ある $e \geq 0$ に対し
$\text{Trace}_{L/K}(\alpha^e)$ が非零であることを示さなければならない。
$P = x^d + a_1 x^{d - 1} + \ldots + a_d$ を $K$ 上の $\alpha$ の
最小多項式とする。このとき補題
\ref{fields-lemma-characteristic-vs-minimal-polynomial} により、$P$ は線形写像
$\alpha : L \to L$ の特性多項式でもある。
% P01-E0255
$L/k$ は分離的なので、補題 \ref{fields-lemma-recognize-separable} により $P$ は
$K$ の代数閉包 $\overline{K}$ に相異なる $d$ 個の根
$\alpha_1, \ldots, \alpha_d$ をもつ。したがって、これらは
$\alpha : L \to L$ の固有値である。線形代数により、$\alpha^e$ のトレースは
$\alpha_1^e + \ldots + \alpha_d^e$ に等しい。よって補題
\ref{fields-lemma-sums-of-powers} から結論が従う。
\end{proof}

\noindent
$K$ を体とし、$Q : V \times V \to K$ を $K$ 上の有限次元ベクトル空間上の
双線形形式とする。$\dim_K(V) = n$ とする。このとき $Q$ は線形写像
$Q : V \to V^*$、$v \mapsto Q(v, -)$ を定める。ただし
$V^* = \Hom_K(V, K)$ は双対ベクトル空間である。したがって線形写像
$$
\det(Q) : \wedge^n(V) \longrightarrow \wedge^n(V)^*
$$
を得る。基底元 $\omega \in \wedge^n(V)$ を選ぶと、
$\det(Q)(\omega) = \lambda \omega^*$ と書ける。ただし $\omega^*$ は
$\wedge^n(V)^*$ の双対基底元である。$\omega$ の選択を、ある $c \in K^*$ に対する
$c \omega$ に変更すると、$\omega^*$ は $c^{-1} \omega^*$ に変わるので、
$\lambda$ は $c^2 \lambda$ に変わる。したがって $K/(K^*)^2$ における
$\lambda$ の類はよく定まり、これを {\it $Q$ の判別式}という。定義を展開すると、
$$
\lambda = \det(Q(v_i, v_j)_{1 \leq i, j \leq n})
$$
である。ただし $\{v_1, \ldots, v_n\}$ は $K$ 上の $V$ の基底である。
判別式が非零であることと $Q$ が非退化であることは同値である。

\begin{definition}
\label{fields-definition-discriminant}
$L/K$ を体の有限次拡大とする。{\it $L/K$ の判別式}とは、
トレース形式 $Q_{L/K}$ の判別式のことである。
\end{definition}

\noindent
上の議論と補題 \ref{fields-lemma-separable-trace-pairing} により、判別式が非零であることと
$L/K$ が分離的であることは同値である。$a \in K$ に対して、より正確には
判別式が $K/(K^*)^2$ における $a$ の類であると言うべきところを、しばしば
「判別式は $a$ である」と言う。

\begin{exercise}
\label{fields-exercise-quadratic-discriminant}
$L/K$ を次数 $2$ の拡大とする。次のうちちょうど一つが起こることを示せ。
\begin{enumerate}
\item 判別式は $0$、$K$ の標数は $2$ であり、$L/K$ は $K$ の元の平方根を
とることによって得られる純非分離拡大である。
\item 判別式は $1$、$K$ の標数は $2$ であり、$L/K$ は次数 $2$ の分離拡大である。
\item 判別式は平方ではなく、$K$ の標数は $2$ ではなく、$L$ は $K$ に
判別式の平方根を添加して得られる。
\end{enumerate}
\end{exercise}






\section{ガロア理論}
\label{fields-section-galois-theory}

\noindent
まず定義を与える。

\begin{definition}
\label{fields-definition-galois}
体拡大 $E/F$ が代数的、分離的、かつ正規であるとき、これを {\it ガロア拡大}
という。
\end{definition}

\noindent
有限次拡大がガロア拡大であることと、それが「正しい」個数の自己同型をもつことは
同値である。

\begin{lemma}
\label{fields-lemma-finite-Galois}
$E/F$ を体の有限次拡大とする。このとき、$E$ が $F$ 上ガロアであることと
$|\text{Aut}(E/F)| = [E : F]$ は同値である。
\end{lemma}

\begin{proof}
$|\text{Aut}(E/F)| = [E : F]$ と仮定する。補題
\ref{fields-lemma-normal-and-automorphisms} により、これは $E/F$ が分離的かつ正規で
あることを含意するので、$E/F$ はガロア拡大である。逆に、$E/F$ が分離的ならば
$[E : F] = [E : F]_s$ であり、さらに正規ならば、補題
\ref{fields-lemma-normal-and-automorphisms} により
$|\text{Aut}(E/F)| = [E : F]$ である。
\end{proof}

\noindent
上の補題に動機づけられ、ガロア群を次のように導入する。

\begin{definition}
\label{fields-definition-galois-group}
$E/F$ がガロア拡大ならば、群 $\text{Aut}(E/F)$ を {\it ガロア群}といい、
$\text{Gal}(E/F)$ と書く。
\end{definition}

\noindent
$L/K$ が無限次ガロア拡大ならば、そのガロア群を位相群とみなすべきである。
これについては Section \ref{fields-section-infinite-galois} で再び扱う。

\begin{lemma}
\label{fields-lemma-galois-goes-up}
$K/E/F$ を代数的体拡大の塔とする。$K$ が $F$ 上ガロアならば、$K$ は
$E$ 上ガロアである。
\end{lemma}

\begin{proof}
補題 \ref{fields-lemma-normal-goes-up} と補題
\ref{fields-lemma-separable-goes-up} を組み合わせればよい。
\end{proof}

\begin{lemma}
\label{fields-lemma-normal-closure-galois}
$L/K$ を体の有限次分離拡大とする。$M$ を $K$ 上の $L$ の正規閉包とする
（定義 \ref{fields-definition-normal-closure}）。このとき $M/K$ はガロア拡大である。
\end{lemma}

\begin{proof}
補題 \ref{fields-lemma-separable-first} の部分拡大 $M/M_{sep}/K$ は、補題
\ref{fields-lemma-separable-first-normal} により正規である。$L/K$ は分離的なので
$L \subset M_{sep}$ である。最小性により $M = M_{sep}$ となり、証明が完了する。
\end{proof}

\noindent
$G$ を体 $K$ に（体の自己同型によって）作用する群とする。しばしば
$$
K^G = \{x \in K \mid \sigma(x) = x \ \forall \sigma \in G\}
$$
と書き、これを $K$ への $G$ の作用の {\it 固定体}という。

\begin{lemma}
\label{fields-lemma-galois-over-fixed-field}
$K$ を体とし、$G$ を $K$ に忠実に作用する有限群とする。このとき拡大
$K/K^G$ はガロアであり、$[K : K^G] = |G|$ で、そのガロア群は $G$ である。
\end{lemma}

\begin{proof}
$\alpha \in K$ が与えられたとき、群作用による $\alpha$ の軌道
$G \cdot \alpha \subset K$ を考える。多項式
$$
P = \prod\nolimits_{\beta \in G \cdot \alpha} (x - \beta) \in K[x]
$$
を考える。補題全体の要点は、この多項式が $G$ の作用で不変であり、したがって
係数が $K^G$ に属することである。実際、$\tau \in G$ に対して
$$
P^\tau = \prod\nolimits_{\beta \in G \cdot \alpha} (x - \tau(\beta)) =
\prod\nolimits_{\beta \in G \cdot \alpha} (x - \beta) = P
$$
である。写像 $\beta \mapsto \tau(\beta)$ が軌道 $G \cdot \alpha$ の置換だからである。
したがって $P \in K^G[x]$ である。また、$\alpha$ は自分の軌道の元なので
$P(\alpha) = 0$ でもある。よって拡大 $K/K^G$ は代数的である。
さらに、$K^G$ 上の $\alpha$ の最小多項式 $Q$ は、今構成した多項式 $P$ を割る。
したがって $Q$ は分離的であり（例えば補題
\ref{fields-lemma-recognize-separable} による）、$K/K^G$ は分離的であると結論する。
よって $K/K^G$ はガロア拡大である。証明を完了するには
$[K : K^G] = |G|$ を示せば十分である。そうすれば補題
\ref{fields-lemma-finite-Galois} により $G$ がガロア群となるからである。

\medskip\noindent
有限個の元 $\alpha_i \in K$、$i = 1, \ldots, n$ を、
$i = 1, \ldots, n$ のすべてについて $\sigma(\alpha_i) = \alpha_i$ ならば
$\sigma$ が $G$ の単位元となるように選ぶ。
$$
L = K^G(\{\sigma(\alpha_i); 1 \leq i \leq n, \sigma \in G\}) \subset K
$$
とおき、$K$ への $G$ の作用が $L$ への $G$ の作用を誘導することに注意する。
$L$ の $K^G$ 上の次数が $|G|$ であることを示す。これで証明は完了する。
実際、$L \subset K$ が真ならば、元 $\alpha \in K$、$\alpha \not \in L$ を
$\alpha_1, \ldots, \alpha_n$ のリストに加えても $L$ を大きくしないことになり、
矛盾する。したがって $K/K^G$ が有限である場合に帰着され、これは次の段落で扱う。

\medskip\noindent
$K/K^G$ が有限であると仮定する。補題
\ref{fields-lemma-primitive-element} により、$K = K^G(\alpha)$ となる
$\alpha \in K$ を見つけられる。この証明の第1段落の構成により、
% P01-E0256
$\alpha$ の $K$ 上の次数は高々 $|G|$ である。一方、構成により $G$ は
$K^G(\alpha) = L$ に忠実に作用し、補題
\ref{fields-lemma-normal-and-automorphisms} の不等式が成り立つので、次数が $|G|$ より
小さくなることはない。
\end{proof}

\begin{theorem}[ガロア理論の基本定理]
\label{fields-theorem-galois-theory}
$L/K$ をガロア群 $G$ をもつ有限次ガロア拡大とする。このとき $K = L^G$ であり、写像
$$
\{\text{subgroups of }G\}
\longrightarrow
\{\text{subextensions }L/M/K\},\quad
H \longmapsto L^H
$$
は全単射で、その逆写像は $M$ を $\text{Gal}(L/M)$ に送る。$G$ の正規部分群
$H$ は、$M/K$ がガロアとなる部分拡大 $M$ とちょうど対応する。
\end{theorem}

\begin{proof}
補題 \ref{fields-lemma-galois-goes-up} により、部分拡大 $L/M/K$ が与えられると、
拡大 $L/M$ はガロアである。もちろん $L/M$ は有限でもある（補題
\ref{fields-lemma-finite-goes-up}）。したがって補題
\ref{fields-lemma-finite-Galois} により
$|\text{Gal}(L/M)| = [L : M]$ である。逆に $H \subset G$ が有限部分群ならば、
補題 \ref{fields-lemma-galois-over-fixed-field} により $[L : L^H] = |H|$ である。
この二つの観察から、定理に述べた全単射な対応が形式的に得られる。

\medskip\noindent
$H \subset G$ が正規ならば、$L^H$ は $G$ の作用で保たれ、標準的な写像
$G/H \to \text{Aut}(L^H/K)$ を得る。この写像は
$\text{Gal}(L/L^H) = H$ なので単射でなければならない。したがって
$|G/H| = [L^H : K]$ であり、補題
\ref{fields-lemma-finite-Galois} により $L^H$ はガロアである。

\medskip\noindent
逆に、$K \subset M \subset L$ で $M/K$ がガロアであると仮定する。補題
\ref{fields-lemma-lift-maps} により、各元 $\tau \in \text{Gal}(L/K)$ は
$\tau|_M \in \text{Gal}(M/K)$ を誘導する。これによりガロア群の準同型
$\text{Gal}(L/K) \to \text{Gal}(M/K)$ が誘導され、その核は $H$ である。
したがって $H$ は正規部分群である。
\end{proof}

\begin{lemma}
\label{fields-lemma-ses-galois}
$L/M/K$ を体の塔とする。$L/K$ と $M/K$ が有限次ガロア拡大であると仮定する。
このとき有限群の短完全列
$$
1 \to \text{Gal}(L/M) \to \text{Gal}(L/K) \to \text{Gal}(M/K) \to 1
$$
を得る。
\end{lemma}

\begin{proof}
実際、補題 \ref{fields-lemma-lift-maps} により、各元
$\tau \in \text{Gal}(L/K)$ は元 $\tau|_M \in \text{Gal}(M/K)$ を誘導し、
これが右側の準同型を与える。左側の写像は、左の群を右向き矢印の核と同一視する。
有限次拡大の塔における次数の乗法性（補題
\ref{fields-lemma-multiplicativity-degrees}）により群の位数が整合するので、列は完全である。
補題 \ref{fields-lemma-lift-maps} を直接用いて、右側の写像が全射であることを
確かめることもできる。
\end{proof}

\section{無限ガロア理論}
\label{fields-section-infinite-galois}

\noindent
ガロア群には標準的な位相が備わる。

\begin{lemma}
\label{fields-lemma-galois-profinite}
$E/F$ をガロア拡大とする。$E$ に離散位相を入れたときに
$$
\text{Gal}(E/F) \times E \longrightarrow E
$$
が連続となるような最弱の位相を $\text{Gal}(E/F)$ に入れる。このとき
\begin{enumerate}
\item 任意の位相空間 $X$ と写像 $X \to \text{Gal}(E/F)$ に対し、
作用 $X \times E \to E$ が連続ならば、誘導される写像
$X \to \text{Gal}(E/F)$ は連続であり、
\item この位相により $\text{Gal}(E/F)$ は副有限位相群となる。
\end{enumerate}
\end{lemma}

\begin{proof}
この証明を通じて $E$ を離散位相空間とみなす。自己写像全体の集合
$\text{Map}(E, E)$ 上のコンパクト開位相は、作用
$\text{Map}(E, E) \times E \to E$ が連続となる普遍的な位相であった。
正確な主張については Topology, Example
\ref{topology-example-automorphisms-of-a-set} を参照されたい。補題で
$\text{Gal}(E/F)$ に入れた位相は、単射
$\text{Gal}(E/F) \to \text{Map}(E, E)$ から誘導される位相である。
したがって普遍性 (1) は、コンパクト開位相の対応する普遍性から従う。
可逆な自己写像全体の集合 $\text{Aut}(E)$ は、コンパクト開位相を入れると
位相群をなす（Topology, Example
\ref{topology-example-automorphisms-of-a-set} を参照）。また
$\text{Gal}(E/F) = \text{Aut}(E/F) \to \text{Map}(E, E)$ は
$\text{Aut}(E)$ を経由するので、位相群を得る。言い換えれば、単射
$$
\text{Gal}(E/F) \subset \text{Aut}(E)
$$
を用いて $\text{Gal}(E/F)$ に誘導位相群構造を入れている
（Topology, Section \ref{topology-section-topological-groups} を参照）。
構成から、これは作用 $\text{Gal}(E/F) \times E \to E$ が連続となる
最弱の位相群構造である。

\medskip\noindent
$\text{Gal}(E/F)$ が副有限であることを示すため、次のように論じる
（ガロア群が有限群の逆極限であることを示す前に位相を定義したので、
この議論は必然的に標準的なものとは異なる）。
Topology, Lemma \ref{topology-lemma-profinite-group} により、
$\text{Gal}(E/F)$ の台となる位相空間が副有限であることを示せば十分である。
任意の部分集合 $S \subset E$ に対して、集合
$$
G(S) = \{ f : S \to E \mid
\begin{matrix}
f(\alpha)\text{ is a root of the minimal polynomial}\\
\text{of }\alpha\text{ over }F\text{ for all }\alpha \in S
\end{matrix}
\}
$$
を考える。多項式の根は有限個しかないので、有限なすべての
$S \subset E$ に対して $G(S)$ は有限である。$S \subset S'$ ならば、
制限により写像 $G(S') \to G(S)$ を得る。また
$\alpha \in S \cap F$ かつ $f \in G(S)$ ならば、この場合の最小多項式は
一次なので $f(\alpha) = \alpha$ であることにも注意する。副有限位相空間
$$
G = \lim_{S \subset E\text{ finite}} G(S)
$$
を考える。標準的な写像
$$
c : \text{Gal}(E/F) \longrightarrow G,\quad
\sigma \longmapsto (\sigma|_S : S \to E)_S
$$
を考える。これは単射であり、定義を展開すれば、上で定義した
$\text{Gal}(E/F)$ の位相は $G$ からの誘導位相であることが分かる。
元 $(f_S) \in G$ が $c$ の像に属するための必要十分条件は、意味をもつとき
（すなわち $\alpha, \beta, \alpha + \beta, \alpha\beta \in S$ のとき）、
(A) $f_S(\alpha) + f_S(\beta) = f_S(\alpha + \beta)$ および
(M) $f_S(\alpha)f_S(\beta) = f_S(\alpha\beta)$ が成り立つことである。
実際、この条件は $\lim f_S : E \to E$ が $F$-代数写像となることを意味し、
したがって補題 \ref{fields-lemma-algebraic-extension-self-map} により自己同型となる。
与えられた三つ組 $(S, \alpha, \beta)$ に対する条件 (A) と (M) は
$G$ の閉部分集合を定める。ゆえに $\text{Gal}(E/F)$ は副有限空間の
閉部分集合と同相であり、それ自身も副有限である。
\end{proof}

\begin{lemma}
\label{fields-lemma-galois-infinite}
$L/M/K$ を体の塔とする。$L/K$ と $M/K$ がともにガロアであると仮定する。
このとき標準的な全射連続準同型
$c : \text{Gal}(L/K) \to \text{Gal}(M/K)$ が存在する。
\end{lemma}

\begin{proof}
補題 \ref{fields-lemma-lift-maps} により、$\text{Gal}(L/K)$ の元
$\tau : L \to L$ が与えられると、その制限 $\tau|_M : M \to M$ は
$\text{Gal}(M/K)$ の元である。これが準同型 $c$ を定める。
連続性は位相の普遍性から従う。実際、$M \subset L$ であり、作用
$\text{Gal}(L/K) \times L \to L$ が連続なので、作用
$$
\text{Gal}(L/K) \times M \longrightarrow M,\quad
(\tau, x) \longmapsto \tau(x) = c(\tau)(x)
$$
は連続である。したがって補題 \ref{fields-lemma-galois-profinite} の (1) により
$c$ は連続である。補題 \ref{fields-lemma-lift-maps} は、この写像が全射であることも
示している。
\end{proof}

\noindent
無限次ガロア拡大のガロア群を考える、より標準的な方法を述べる。

\begin{lemma}
\label{fields-lemma-infinite-galois-limit}
$L/K$ をガロア群 $G$ をもつガロア拡大とする。$\Lambda$ を有限次ガロア
部分拡大全体の集合とする。すなわち、$\lambda \in \Lambda$ は
$L/L_\lambda/K$ に対応し、$L_\lambda/K$ はガロア群
$G_\lambda$ をもつ有限次ガロア拡大である。$\lambda \geq \lambda'$ であることを
$L_\lambda \supset L_{\lambda'}$ であることと定めて、$\Lambda$ に半順序を入れる。
このとき
\begin{enumerate}
\item $\Lambda$ は有向半順序集合であり、
\item $L_\lambda$ は $\Lambda$ 上の $K$-拡大の系であって
$L = \colim L_\lambda$ であり、
\item $G_\lambda$ は $\Lambda$ 上の有限群の逆系で、遷移写像は全射であり、
$$
G = \lim_{\lambda \in \Lambda} G_\lambda
$$
が副有限群として成り立ち、
\item 各射影 $G \to G_\lambda$ は連続かつ全射である。
\end{enumerate}
\end{lemma}

\begin{proof}
$K$ を含む $L$ のすべての部分体は $K$ 上分離的である
（これは定義から直ちに従う）。$S \subset L$ を有限部分集合とする。
このとき $K(S)/K$ は有限次であり、$E/K$ が有限次ガロアとなるような塔
$L/E/K(S)/K$ が存在する。補題
\ref{fields-lemma-normal-closure-inside-normal} を参照されたい。
したがってある $\lambda \in \Lambda$ に対して $E = L_\lambda$ である。
特に、集合 $\Lambda$ は空でない。また
$\lambda_1, \lambda_2 \in \Lambda$ が与えられると、有限集合
$S_1, S_2 \subset L$ を用いて $L_{\lambda_i} = K(S_i)$ と書ける
（補題 \ref{fields-lemma-finite-finitely-generated}）。このとき
$K(S_1 \cup S_2) \subset L_\lambda$ となる $\lambda \in \Lambda$ が存在する。
ゆえに $\lambda \geq \lambda_1, \lambda_2$ であり、$\Lambda$ は有向である
（Categories, Definition \ref{categories-definition-directed-system}）。
最後に、$L$ の各元はある $\lambda \in \Lambda$ に対する
$L_\lambda$ に含まれるので、Categories, Section
\ref{categories-section-directed-colimits} のフィルター余極限の記述から
$\colim L_\lambda = L$ が従う。

\medskip\noindent
$\Lambda$ において $\lambda \geq \lambda'$ ならば、補題
\ref{fields-lemma-ses-galois} により標準的な全射
$G_\lambda \to G_{\lambda'}$、$\sigma \mapsto \sigma|_{L_{\lambda'}}$ を得る。
したがって、全射な遷移写像をもつ有限群の逆系を得る。

\medskip\noindent
$G = \text{Aut}(L/K)$ であることを思い出す。補題
\ref{fields-lemma-galois-infinite} により、$\sigma \in G$ の $L_\lambda$ への制限
$\sigma|_{L_\lambda}$ は $G_\lambda$ の元である。さらに、この手続きは連続な全射
$G \to G_\lambda$ を与える。$G_\lambda$ の逆系における遷移写像も制限で
与えられるので、標準的な連続写像
$$
G \longrightarrow \lim_{\lambda \in \Lambda} G_\lambda
$$
を得ることは明らかである。連続性は位相群の圏における極限の定義から従う。
これらの極限が集合および位相空間の圏への忘却関手と可換であることを
Topology, Lemma \ref{topology-lemma-topological-group-limits} から思い出そう。
一方、$L = \colim L_\lambda$ なので、逆極限の任意の元を集合として見れば
$L$ の自己同型を定めることは明らかである。したがってこの写像は全単射である。
両辺の位相は副有限であり、副有限空間の間の全単射連続写像は同相写像である
（Topology, Lemma \ref{topology-lemma-bijective-map}）から、証明が完了する。
\end{proof}

\begin{theorem}[無限ガロア理論の基本定理]
% P01-E0257
\label{fields-theorem-inifinite-galois-theory}
$L/K$ をガロア拡大とする。$G = \text{Gal}(L/K)$ を、副有限位相群とみなした
ガロア群とする（補題 \ref{fields-lemma-galois-profinite}）。このとき $K = L^G$ であり、
写像
$$
\{\text{closed subgroups of }G\}
\longrightarrow
\{\text{subextensions }L/M/K\},\quad
H \longmapsto L^H
$$
は全単射で、その逆写像は $M$ を $\text{Gal}(L/M)$ に送る。
有限次部分拡大 $M$ は、開部分群 $H \subset G$ とちょうど対応する。
$G$ の正規閉部分群 $H$ は、$K$ 上ガロアとなる部分拡大 $M$ とちょうど対応する。
\end{theorem}

\begin{proof}
有限ガロア理論の結果（定理 \ref{fields-theorem-galois-theory}）を断りなく用いる。
$S \subset L$ を有限部分集合とする。$K(S) \subset E$ かつ
$E/K$ が有限次ガロアとなるような塔 $L/E/K$ が存在する。補題
\ref{fields-lemma-normal-closure-inside-normal} を参照されたい。言い換えれば、
$L/K$ はその有限次ガロア部分拡大の合併である。こうした $E$ に対し、
補題 \ref{fields-lemma-galois-infinite} により写像
$\text{Gal}(L/K) \to \text{Gal}(E/K)$ は全射かつ連続である。すなわち、
$\text{Gal}(E/K)$ の位相は離散なので、その核は開である。特に
$E^{\text{Gal}(E/K)} = K$ だから、$L \setminus K$ のどの元も
$\text{Gal}(L/K)$ によって固定されないことが分かる。これで $L^G = K$ が示された。

\medskip\noindent
部分拡大 $L/M/K$ が与えられると、補題 \ref{fields-lemma-galois-goes-up} により
拡大 $L/M$ はガロアである。$G$ 上の位相の定義から、部分群
$\text{Gal}(L/M)$ が閉であることは直ちに分かる。上の議論を $L/M$ に適用すれば
$L^{\text{Gal}(L/M)} = M$ を得る。

\medskip\noindent
逆に、$H \subset G$ を閉部分群とする。$H = \text{Gal}(L/L^H)$ であることを
示す。包含 $H \subset \text{Gal}(L/L^H)$ は明らかである。
$g \in \text{Gal}(L/L^H)$ と仮定し、$S \subset L$ を有限部分集合とする。
% P01-E1270
$g$ の開近傍 $U_S(g) = \{g' \in G \mid g'(s) = g(s)\}$ が $H$ と交わることを
示す。$H$ は閉なので、これは $g \in H$ を意味する。証明の第1段落と同様に、
$K(S)$ を含む有限次ガロア部分拡大 $L/E/K$ を取り、準同型
$c : \text{Gal}(L/K) \to \text{Gal}(E/K)$ を考える。このとき
$L^H \cap E = E^{c(H)}$ である。$g$ は $L^H$ を固定するので
$E^{c(H)}$ を固定し、有限ガロア理論により $c(g) \in c(H)$ である。
$c(h) = c(g)$ となる $h \in H$ を取る。このとき望みどおり
$h \in U_S(g)$ である。

\medskip\noindent
以上で、閉部分群と部分拡大の間の対応が確立された。

\medskip\noindent
$H \subset G$ が開であると仮定する。上と同様に論じると、十分大きな有限次
ガロア部分拡大 $E$ に対して $H$ は $\text{Gal}(L/E)$ を含む。したがって
$L^H$ は $E$ に含まれ、ゆえに $K$ 上有限次である。逆に $M$ が有限次部分拡大
ならば、$M$ は有限部分集合 $S$ で生成され、対応する部分群は開集合
$U_S(e)$ である。ここで $e \in G$ は単位元である。

\medskip\noindent
$K \subset M \subset L$ かつ $M/K$ がガロアであると仮定する。
補題 \ref{fields-lemma-galois-infinite} により、核が $\text{Gal}(L/M)$ である
ガロア群の全射連続準同型
$\text{Gal}(L/K) \to \text{Gal}(M/K)$ が存在する。したがって
$\text{Gal}(L/M)$ は正規閉部分群である。

\medskip\noindent
最後に、$N \subset G$ を正規閉部分群とする。証明の第1段落に現れる任意の
$L/E/K$ に対して、像 $c(N) \subset \text{Gal}(E/K)$ は正規部分群である。
したがって $L^N = \bigcup E^{c(N)}$ は $K$ のガロア拡大の合併であり
（有限ガロア理論による）、ゆえに $K$ 上ガロアである。
\end{proof}

\begin{lemma}
\label{fields-lemma-ses-infinite-galois}
$L/M/K$ を体の塔とする。$L/K$ と $M/K$ がガロアであると仮定する。
このとき副有限位相群の短完全列
$$
1 \to \text{Gal}(L/M) \to \text{Gal}(L/K) \to \text{Gal}(M/K) \to 1
$$
を得る。
\end{lemma}

\begin{proof}
これは補題 \ref{fields-lemma-galois-infinite} の言い換えである。
\end{proof}

\section{複素数}
\label{fields-section-complex-numbers}

\noindent
代数学の基本定理は、複素数体が代数閉体であると主張する。この節では
このことを簡単に論じる。

\medskip\noindent
最初に注意したいのは、これを証明するには解析学から多少の結果を用いる
必要があるということである。実係数の任意の奇数次多項式は実根をもつという、
直観的に明らかな事実を用いる。実際、ある $k \geq 0$ と
$a_{2k + 1} \not = 0$ に対して
$P(x) = a_{2k + 1} x^{2k + 1} + \ldots + a_0 \in \mathbf{R}[x]$
とする。$a_{2k + 1} > 0$ と仮定してよいし、そう仮定する。
$x \in \mathbf{R}$ が十分大きな正数ならば、項
$a_{2k + 1} x^{2k + 1}$ が他のすべての項を支配するので $P(x) > 0$ である。
同様に、$x \ll 0$ ならば同じ理由により $P(x) < 0$ である
（ここで次数が奇数であることを用いる）。したがって中間値の定理により、
$P(x) = 0$ となる $x \in \mathbf{R}$ が存在する。

\medskip\noindent
以上から、$\mathbf{R}$ は自明でない奇数次の体拡大をもたないことが分かる。
実際、そのような拡大の元は奇数次の最小多項式をもつことになる。

\medskip\noindent
次に、$K/\mathbf{R}$ をガロア群 $G$ をもつ有限次ガロア拡大とする。
$P \subset G$ を $2$-Sylow 部分群とする。このとき
$K^P/\mathbf{R}$ は奇数次拡大なので、上の結果により
$K^P = \mathbf{R}$ であり、ひいては $G = P$ である
（もちろん、これらの議論はすべてガロア理論に依存している）。
したがって $G$ は $2$-群である。$G$ が自明でなければ、
$\mathbf{C}$ は同型を除いて $\mathbf{R}$ の唯一の次数 $2$ の拡大なので
$\mathbf{C} \subset K$ である。$G$ が $2$ 個より多くの元をもつならば、
$\mathbf{C}$ の二次拡大を得ることになる。しかし、すべての複素数は
平方根をもつので、これは矛盾である。

\medskip\noindent
結論として、$\mathbf{C}$ は代数閉である。実際、そうでなければ
自明でない有限次拡大 $K/\mathbf{C}$ を得る。補題
\ref{fields-lemma-normal-closure} により、これを $\mathbf{R}$ 上正規
（したがってガロア）であると仮定してよい。しかし上で見たように、
その場合 $K = \mathbf{C}$ である。

\begin{lemma}[代数学の基本定理]
\label{fields-lemma-C-algebraically-closed}
体 $\mathbf{C}$ は代数閉である。
\end{lemma}

\begin{proof}
上の議論を参照されたい。
\end{proof}

\section{クンマー拡大}
\label{fields-section-Kummer}

\noindent
$K$ を体とする。$n \geq 2$ を、$K$ が $1$ の原始 $n$ 乗根を含むような
整数とする。% P01-E1271
$a \in K$ とする。$L$ を、方程式 $x^n = a$ の根 $b$ を添加して得られる
$K$ の拡大とする。このとき $L/K$ はガロアである。
$G = \text{Gal}(L/K)$ をガロア群とすると、写像
$$
G \longrightarrow \mu_n(K),\quad \sigma \longmapsto \sigma(b)/b
$$
は群の単射準同型である。特に、$G$ は巡回群 $\mu_n(K)$ の部分群として、
位数が $n$ を割る巡回群である。クンマー理論はこの逆を与える。

\begin{lemma}[クンマー拡大]
\label{fields-lemma-Kummer}
$L/K$ を、ガロア群が $\mathbf{Z}/n\mathbf{Z}$ である体のガロア拡大とする。
さらに $K$ の標数は $n$ と互いに素であり、$K$ は $1$ の原始 $n$ 乗根を
含むと仮定する。このとき $L = K[z]$ かつ $z^n \in K$ である。
\end{lemma}

\begin{proof}
$\zeta \in K$ を $1$ の原始 $n$ 乗根とする。
$\sigma$ を $\text{Gal}(L/K)$ の生成元とする。
$\sigma : L \to L$ を $K$-線形作用素とみなす。線形作用素として
$\sigma^n - 1 = 0$ であることに注意する。指標の線形独立性
（補題 \ref{fields-lemma-independence-characters}）を適用すると、
$\sigma$ を零化する次数 $< n$ の $K$ 上の多項式は存在しないことが分かる。
したがって、線形作用素としての $\sigma$ の最小多項式は $x^n - 1$ である。
$\zeta$ は $x^n - 1$ の根なので、線形代数により
$\sigma(z) = \zeta z$ となる $0 \neq z \in L$ が存在する。この $z$ は
$\sigma(z^n) = (\zeta z)^n = z^n$ なので $z^n \in K$ を満たす。
さらに
$z, \sigma(z), \ldots, \sigma^{n - 1}(z) =
z, \zeta z, \ldots \zeta^{n - 1} z$ は相異なるので、$z$ は $K$ 上 $L$ を
生成することが保証される。ゆえに求めるとおり $L = K[z]$ である。
\end{proof}

\begin{lemma}
\label{fields-lemma-adjoint-pth-root-unity}
$K$ を代数閉包 $\overline{K}$ をもつ体とする。
$p$ を $K$ の標数とは異なる素数とする。
$\zeta \in \overline{K}$ を $1$ の原始 $p$ 乗根とする。このとき
$K(\zeta)/K$ は次数が $p - 1$ を割るガロア拡大である。
\end{lemma}

\begin{proof}
多項式 $x^p - 1$ の根は
$1, \zeta, \zeta^2, \ldots, \zeta^{p - 1}$ なので、
$K(\zeta)$ 上で完全に分解する。したがって $K(\zeta)/K$ は分解体であり、
ゆえに正規である。$x^p - 1$ は分離多項式なので、この拡大は分離的である。
したがって拡大はガロアである。$K$ 上の $K(\zeta)$ の任意の自己同型は、
ある $1 \leq i \leq p - 1$ に対して $\zeta$ を $\zeta^i$ に送る。
ゆえにガロア群は $(\mathbf{Z}/p\mathbf{Z})^*$ の部分群である。
\end{proof}

\begin{lemma}
\label{fields-lemma-subfields-kummer}
\begin{reference}
\cite[Theorem 5.2]{Radical}
\end{reference}
$K$ を体とする。$L/K$ を次数 $e$ の有限次拡大で、
$a = \alpha^e \in K$ となる元 $\alpha$ で生成されるものとする。
$L$ に属するすべての $e$ 乗根が $K$ に含まれるならば、
% P01-E0258
任意の部分拡大 $L/L'/K$ は、ある $d | e$ に対して $\alpha^d$ で生成される。
\end{lemma}

\begin{proof}
$d | e$ に対して部分体 $K(\alpha^d)$ は
$[K(\alpha^d) : K] = e/d$ および $[L : K(\alpha^d)] = d$ を満たすことに
注意する。$L/L'/K$ を部分拡大とし、$d = [L : L']$ とおく。
$\alpha^d \in L'$ ならば、次数を比較して $L' = K(\alpha^d)$ である。
$P \in L'[x]$ を $L'$ 上の $\alpha$ の最小多項式とする。
このとき $P$ は $x^e - a$ を割り、$P$ の次数は $d$ である。
$x^e - a$ の $L$ 上の分解体において
$$
x^e - a = \prod\nolimits_{i = 1, \ldots, e} (x - \zeta_i \alpha)
$$
と書く。$\zeta_i$ は $e$ 乗根であり、番号を付け替えれば
$$
P = \prod\nolimits_{i = 1, \ldots, d} (x - \zeta_i \alpha)
$$
となる。$P$ の定数項は
% P01-E0259
$$
c = (\prod\nolimits_{i = 1, \ldots, d} \zeta_i) \alpha^d
$$
に等しく、$L' \subset L$ に属する。$\alpha \in L$ なので、
$\zeta = \prod_{i = 1, \ldots, d} \zeta_i$ は $L$ に属し、仮定により
$K$ に属する。したがって $\alpha^d = \zeta^{-1}c \in L'$ となり、
結論を得る。
\end{proof}

\section{アルティン--シュライアー拡大}
\label{fields-section-Artin-Schreier}

\noindent
$K$ を標数 $p > 0$ の体とする。$a \in K$ とする。$L$ を、方程式
$x^p - x = a$ の根 $b$ を添加して得られる $K$ の拡大とする。
このとき $L/K$ はガロアである。$G = \text{Gal}(L/K)$ をガロア群とすると、
写像
$$
G \longrightarrow \mathbf{Z}/p\mathbf{Z},\quad
\sigma \longmapsto \sigma(b) - b
$$
は群の単射準同型である。特に、$G$ は $\mathbf{Z}/p\mathbf{Z}$ の部分群として、
位数が $p$ を割る巡回群である。アルティン--シュライアー拡大の理論は
この逆を与える。

\begin{lemma}[アルティン--シュライアー拡大]
\label{fields-lemma-Artin-Schreier}
$L/K$ を、ガロア群が $\mathbf{Z}/p\mathbf{Z}$ である標数 $p > 0$ の
体のガロア拡大とする。このとき $L = K[z]$ かつ $z^p - z \in K$ である。
\end{lemma}

\begin{proof}
$\sigma$ を $\text{Gal}(L/K)$ の生成元とする。
$\sigma : L \to L$ を $K$-線形作用素とみなす。線形作用素として
$\sigma^p - 1 = 0$ であることに注意する。指標の線形独立性
（補題 \ref{fields-lemma-independence-characters}）を適用すると、
$\sigma$ を零化する次数 $< p$ の多項式は存在しない。
したがって $\sigma$ の最小多項式は $x^p - 1 = (x - 1)^p$ である。
これは、$(\sigma - 1)^{p - 1}(w) = y$ が零でないような
$w \in L$ が存在することを意味する。このとき $\sigma(y) = y$、すなわち
$y \in K$ である。したがって
$z = y^{-1}(\sigma - 1)^{p - 2}(w)$ は $\sigma(z) = z + 1$ を満たす。
$z \not \in K$ なので $L = K[z]$ である。さらに
$\sigma(z^p - z) = (z + 1)^p - (z + 1) = z^p - z$ なので
$z^p - z \in K$ であり、証明が完了する。
\end{proof}

\section{超越性}
\label{fields-section-transcendence}

\noindent
標準的な定義を思い出す。

\begin{definition}
\label{fields-definition-transcendence}
$K/k$ を体拡大とする。
\begin{enumerate}
\item $K$ の元の族 $\{x_i\}_{i \in I}$ が $k$ 上{\it 代数的独立}であるとは、
$X_i$ を $x_i$ に送る写像
$$
k[X_i; i\in I] \longrightarrow K
$$
が単射であることをいう。
\item 多項式環 $k[x_i; i \in I]$ の分数体を $k(x_i; i\in I)$ と書く。
\item $k$ の{\it 純超越拡大}とは、$k$ 上の多項式環の分数体と同型な
任意の体拡大 $K/k$ をいう。
\item $K/k$ の{\it 超越基底}とは、$k$ 上代数的独立で、拡大
$K/k(x_i; i\in I)$ が代数的となるような元の族
$\{x_i\}_{i \in I}$ をいう。
\end{enumerate}
\end{definition}

\begin{example}
\label{fields-example-pi-transcendental}
$\pi$ は有理係数の零でない多項式の根ではないので、体
$\mathbf{Q}(\pi)$ は純超越的である。特に
$\mathbf{Q}(\pi) \cong \mathbf{Q}(x)$ である。
\end{example}

\begin{lemma}
\label{fields-lemma-transcendence-degree}
$E/F$ を体拡大とする。$F$ 上の $E$ の超越基底は存在する。
任意の二つの超越基底は同じ濃度をもつ。
\end{lemma}

\begin{proof}
$A$ を $E$ の代数的独立部分集合とする。$G$ を、$A$ を含み、
$E/F$ を生成する $E$ の部分集合とする。
$A \subset B \subset G$ を満たす超越基底 $B$ を見つけられることを示す。
そのため、$G$ の部分集合で $A$ を含む代数的独立集合全体の族
$\mathcal{B}$ を考え、包含関係で $\mathcal{B}$ に半順序を入れる。
このとき $\mathcal{B}$ は少なくとも元 $A$ を含む。
$\mathcal{B}$ の全順序部分集合 $T$ の元の合併は、$F$ 上 $E$ の
代数的独立部分集合である。実際、代数的従属関係があれば、多項式は有限個の
変数しか含まないので、それは $T$ のいずれかの元ですでに生じていたはずである。
この合併は $A$ を含み、$G$ に含まれる。ツォルンの補題により、極大元
$B \in \mathcal{B}$ が存在する。ここで $E$ は $F(B)$ 上代数的であると主張する。
そうでなければ $F(G) = E$ なので、$F(B)$ 上超越的な元 $f \in G$ が存在する。
すると $B \cup\{f\}$ は代数的独立となり、$B$ の極大性に矛盾する。
ゆえに $B$ が求める超越基底である。

\medskip\noindent
$B$ と $B'$ を二つの超越基底とする。一般性を失わず
$|B'| \leq |B|$ と仮定してよい。証明を二つの場合に分ける。
第一の場合として、$B$ が無限集合であるとする。このとき各
$\alpha \in B'$ に対して、$\alpha$ が $F(B_{\alpha})$ 上代数的となるような
有限集合 $B_{\alpha} \subset B$ が存在する。実際、代数的従属関係は
有限個の不定元しか用いない。そこで
$B^* = \bigcup_{\alpha\in B'} B_{\alpha}$ と定める。構成により
$B^* \subset B$ だが、実際には両集合は等しいと主張する。
等しくないとし、$\beta \in B \setminus B^*$ が存在するとする。
$\beta$ は $F(B')$ 上代数的であり、後者は $F(B^*)$ 上代数的なので、
$\beta$ は $F(B^*)$ 上代数的となり、矛盾する。したがって
$|B| \leq |\bigcup_{\alpha \in B'} B_{\alpha}|$ である。
ここで $B'$ が有限ならば $B$ も有限となるので、$B'$ は無限と仮定してよい。
すると
$$
|B| \leq |\bigcup\nolimits_{\alpha \in B'} B_{\alpha}| = |B'|
$$
である。各 $B_\alpha$ は有限で $B'$ は無限だからである。
したがって無限の場合には $|B| = |B'|$ である。

\medskip\noindent
次に $B$ が有限の場合を考える。この場合 $B'$ も有限なので、
$B = \{\alpha_1, \ldots, \alpha_n\}$ および
$B' = \{\beta_1, \ldots, \beta_m\}$、$m \leq n$ とする。
$m$ に関する帰納法を用いる。$m = 0$ ならば $E/F$ は代数的なので
$B = \emptyset$、したがって $n = 0$ である。$m > 0$ ならば、
$f(\beta_1, \alpha_1, \ldots, \alpha_n) = 0$ で、$x$ が $f$ に現れるような
既約多項式 $f \in F[x, y_1, \ldots, y_n]$ が存在する。
$\beta_1$ は $F$ 上代数的でないので、$f$ はある $y_i$ を含まなければならない。
一般性を失わず、$f$ は $y_1$ を含むと仮定する。
$B^* = \{\beta_1, \alpha_2, \ldots, \alpha_n\}$ とおく。
$B^*$ が $E/F$ の基底であると主張する。この主張を示すため、
$\alpha_1$ は $F(B^*)$ 上代数的なので、代数拡大の塔
$$
E/ F(B^*, \alpha_1) / F(B^*)
$$
を得る。さらに、$B^*$ は（元の重複度も数えれば）$F$ 上代数的独立であると
主張する。そうでなければ、
$g(\beta_1, \alpha_2, \ldots, \alpha_n) = 0$ となる既約多項式
$g\in F[x, y_2, \ldots, y_n]$ が存在する。これは $x$ を含まなければならず、
$\beta_1$ は $F(\alpha_2, \ldots, \alpha_n)$ 上代数的となる。
すると $\alpha_1$ も $F(\alpha_2, \ldots, \alpha_n)$ 上代数的となるが、
これは不可能である。したがって
$\{\alpha_2, \ldots, \alpha_n\}$ と
$\{\beta_2, \ldots, \beta_m\}$ は $F(\beta_1)$ 上の $E$ の基底である。
帰納法により $m = n$ となる。
\end{proof}

\begin{definition}
\label{fields-definition-transcendence-degree}
$K/k$ を体拡大とする。$k$ 上の $K$ の{\it 超越次数}とは、
$k$ 上の $K$ の超越基底の濃度である。これを $\text{trdeg}_k(K)$ と書く。
\end{definition}

\begin{lemma}
\label{fields-lemma-transcendence-degree-tower}
$L/K/k$ を体拡大の塔とする。このとき
$$
\text{trdeg}_k(L) =
\text{trdeg}_K(L) +
\text{trdeg}_k(K).
$$
\end{lemma}

\begin{proof}
$k$ 上の $K$ の超越基底 $A \subset K$ を取る。
$K$ 上の $L$ の超越基底 $B \subset L$ を取る。このとき
$A \cup B$ が $k$ 上の $L$ の超越基底であることは容易に分かる。
\end{proof}

\begin{example}
\label{fields-example-pi-e-transcendental}
数 $e$ と $\pi$ を添加して得られる体拡大 $\mathbf{Q}(e, \pi)$ を考える。
$e$ と $\pi$ はともに有理数体上超越的なので、この体拡大の超越次数は
少なくとも $1$ である。しかし $e$ と $\pi$ が代数的独立ならば、
この体拡大の超越次数は $2$ となる。このことが真かどうかは未知であり、
したがって $\text{trdeg}(\mathbf{Q}(e, \pi))$ を決定する問題は未解決である。
\end{example}

\begin{example}
\label{fields-example-function-field-not-unique-transcendence-basis}
$F$ を体、$E = F(t)$ とする。$E = F(t)$ なので $\{t\}$ は超越基底である。
しかし $F(t)/F(t^2)$ は代数的なので、$\{t^2\}$ も超越基底である。
これは、拡大 $E/F$ を代数拡大 $E/F'$ と純超越拡大 $F'/F$ に常に分解できる
一方、この分解は一意でなく、超越基底の選択に依存することを示している。
\end{example}

\begin{example}
\label{fields-example-riemann-surface-transcendence}
$X$ をコンパクトリーマン面とする。このとき関数体 $\mathbf{C}(X)$
（Example \ref{fields-example-field-of-meromorphic-functions} を参照）は
$\mathbf{C}$ 上超越次数一である。実際、$\mathbf{C}$ の超越次数一の
{\it 任意の}有限生成拡大は、リーマン面から生じる。さらに、コンパクト
リーマン面と非定数正則写像の圏と、$\mathbf{C}$ の超越次数 $1$ の
有限生成拡大と $\mathbf{C}$-代数の射の圏の反対圏との間には圏同値がある。
\cite{Forster} を参照されたい。

\medskip\noindent
上の主張には代数的な版もある。代数閉体 $k$（例えば複素数体）上の射影空間内の
（既約）代数曲線が与えられたとき、その「有理関数体」を考えられる。
これは本質的には、多項式の商の形をした関数で、分母が曲線上恒等的に
零とならないものである。滑らかな射影曲線と曲線の非定数射の圏と、
$k$ の超越次数一の有限生成拡大の圏との間にも同様の反同値がある
（Algebraic Curves, Theorem
\ref{curves-theorem-curves-rational-maps}）。\cite{H} を参照されたい。
\end{example}

\begin{definition}
\label{fields-definition-algebraically-closed-in}
$K/k$ を体拡大とする。
\begin{enumerate}
\item $K$ における{\it $k$ の代数閉包}とは、$k$ 上代数的である
$K$ の元からなる $K$ の部分体 $k'$ をいう。
\item $k$ 上代数的である $K$ のすべての元が $k$ に含まれるとき、
$k$ は{\it $K$ において代数的に閉じている}という。
\end{enumerate}
\end{definition}

\begin{lemma}
\label{fields-lemma-purely-transcendental-degree}
$k'/k$ を体の有限次拡大とする。
$k'(x_1, \ldots, x_r)/k(x_1, \ldots, x_r)$ を、誘導される純超越拡大の拡大
とする。このとき
$[k'(x_1, \ldots, x_r) : k(x_1, \ldots, x_r)] = [k' : k] < \infty$ である。
\end{lemma}

\begin{proof}
拡大次数の乗法性（補題 \ref{fields-lemma-multiplicativity-degrees}）により、
$k'$ が単一の元 $\alpha \in k'$ で $k$ 上生成される場合を証明すれば十分である。
$f \in k[T]$ を $k$ 上の $\alpha$ の最小多項式とする。このとき
$k'(x_1, \ldots, x_r)$ は $k$ 上 $\alpha, x_1, \ldots, x_r$ で生成される。
したがって $k'(x_1, \ldots, x_r)$ は
$k(x_1, \ldots, x_r)$ 上 $\alpha$ で生成される。ゆえに
$f$ が $k(x_1, \ldots, x_r)[T]$ の元としても既約であることを示せば十分である。
証明の概略だけを述べる。$f$ は $k[x_1, \ldots, x_r, T]$ の元として
既約であることは明らかである。例えば、$f$ は $T$ の多項式としてモニックであり、
$k[x_1, \ldots, x_r, T]$ における仮想的な因数分解は、$x_i$ を $0$ とおくことで
$k[T]$ における因数分解を導くからである。ガウスの補題により結論を得る。
\end{proof}

\begin{lemma}
\label{fields-lemma-algebraic-closure-in-finitely-generated}
$K/k$ を有限生成体拡大とする。$K$ における $k$ の代数閉包は $k$ 上有限次である。
\end{lemma}

\begin{proof}
$x_1, \ldots, x_r \in K$ を $k$ 上の $K$ の超越基底とする。このとき
$n = [K : k(x_1, \ldots, x_r)] < \infty$ である。
$k'/k$ が有限次となるような $k \subset k' \subset K$ を仮定する。この場合
$[k'(x_1, \ldots, x_r) : k(x_1, \ldots, x_r)] = [k' : k] < \infty$ である。
補題 \ref{fields-lemma-purely-transcendental-degree} を参照されたい。したがって
$$
[k' : k] = [k'(x_1, \ldots, x_r) : k(x_1, \ldots, x_r)] \leq
[K : k(x_1, \ldots, x_r)] = n.
$$
言い換えれば、有限次部分拡大の次数は有界であり、補題が従う。
\end{proof}

\section{線形無関連拡大}
\label{fields-section-linearly-disjoint}

\noindent
$k$ を体とし、$K$ と $L$ を $k$ の体拡大とする。さらに $K$ と $L$ は
あるより大きな体 $\Omega$ に埋め込まれていると仮定する。

\begin{definition}
\label{fields-definition-compositum}
体拡大の図式
\begin{equation}
\label{fields-equation-inside-omega}
\vcenter{
\xymatrix{
L \ar[r] & \Omega \\
k \ar[r] \ar[u] & K \ar[u]
}
}
\end{equation}
を考える。$\Omega$ における{\it $K$ と $L$ の合成体}とは、
$L$ と $K$ の両方を含む $\Omega$ の最小の部分体であり、$KL$ と書く。
\end{definition}

\noindent
$KL$ は $k$ 上集合 $K \cup L$ で生成され、$L$ 上集合 $K$ で生成され、
$K$ 上集合 $L$ で生成されることは明らかである。

\medskip\noindent
% P01-E0260
{\bf 注意：} 合成体の（同型類）は、$K$ と $L$ の $\Omega$ への埋め込みの
選択に依存する。例えば数体 $K = \mathbf{Q}(2^{1/8}) \subset \mathbf{R}$ と
$L = \mathbf{Q}(2^{1/12}) \subset \mathbf{R}$ を考える。
$\mathbf{R}$ 内の合成体は、$\mathbf{Q}$ 上次数 $24$ の体
$\mathbf{Q}(2^{1/24})$ である。しかし
$K = \mathbf{Q}[x]/(x^8 - 2)$ を、$x$ を $2^{1/8}e^{2\pi i/8}$ に送ることで
$\mathbf{C}$ に埋め込むと、合成体
$\mathbf{Q}(2^{1/12}, 2^{1/8}e^{2\pi i/8})$ は
$i = e^{2\pi i/4}$ を含み、$\mathbf{Q}$ 上次数 $48$ をもつ
（次数が $48$ であることの証明は省略するが、$i$ の存在だけで二つの合成体が
同型でないことは確実に分かる）。

\begin{definition}
\label{fields-definition-linearly-disjoint}
(\ref{fields-equation-inside-omega}) のような体の図式を考える。写像
$$
K \otimes_k L \longrightarrow KL,\quad
\sum x_i \otimes y_i \longmapsto \sum x_i y_i
$$
が単射であるとき、$K$ と $L$ は{\it $\Omega$ において $k$ 上線形無関連}
であるという。
\end{definition}

\noindent
次の補題は、ほかに適切な置き場所がないように思われる。

\begin{lemma}
\label{fields-lemma-normal-case}
$E/F$ を正規代数体拡大とする。次を満たす部分拡大
$E / E_{sep} /F$ および $E / E_{insep} / F$ が存在する。
\begin{enumerate}
\item $F \subset E_{sep}$ はガロアで、$E_{sep} \subset E$ は純非分離的であり、
\item $F \subset E_{insep}$ は純非分離的で、$E_{insep} \subset E$ はガロアであり、
\item $E = E_{sep} \otimes_F E_{insep}$ である。
\end{enumerate}
\end{lemma}

\begin{proof}
部分体 $E_{sep}$ は補題 \ref{fields-lemma-separable-first} で得た。
$E_{insep} = E^{\text{Aut}(E/F)}$ とおく。詳細は省略する。
\end{proof}

\section{復習}
\label{fields-section-algebraic}

\noindent
この節では、上で得られたことを手短に復習する。

\medskip\noindent
$K/k$ を体拡大とし、$\alpha \in K$ とする。このとき次の可能性がある。
\begin{enumerate}
\item 元 $\alpha$ は $k$ 上超越的である。
\item 元 $\alpha$ は $k$ 上代数的である。
% P01-E0261
その{\it 最小多項式}を $P(T) \in k[T]$ と書く。これは $k$ に係数をもつ
モニック多項式 $P(T) = T^d + a_1 T^{d - 1} + \ldots + a_d$ である。
これは既約で $P(\alpha) = 0$ を満たす。これらの性質は $P$ を一意に定め、
整数 $d$ を{\it $k$ 上の $\alpha$ の次数}という。二つの場合に分かれる。
\begin{enumerate}
\item 多項式 $\text{d}P/\text{d}T$ は恒等的に零ではない。
これは、$k$ の代数閉包内の相異なる元
$\alpha_1, \ldots, \alpha_d$ によって
$P(T) = \prod_{i = 1, \ldots, d} (T - \alpha_i)$ と書けることと同値である。
この場合、$\alpha$ は $k$ 上{\it 分離的}であるという。
\item $\text{d}P/\text{d}T$ は恒等的に零である。この場合、
% P01-E0262
$k$ の標数 $p$ は $ > 0$ であり、$P$ は実際に $T^p$ の多項式である。
$P$ が $T^q$ の多項式となる最大の冪 $q = p^e$ が存在することは明らかである。
このとき元 $\alpha^q$ は $k$ 上分離的である。
\end{enumerate}
\end{enumerate}

\begin{definition}
\label{fields-definition-separable-algebraic}
代数体拡大。
\begin{enumerate}
\item $K$ のすべての元が $k$ 上代数的であるとき、体拡大 $K/k$ は
{\it 代数的}であるという。
\item すべての $\alpha \in k'$ が $k$ 上分離的であるとき、代数拡大
$k'/k$ は{\it 分離的}であるという。
\item $k$ の標数が $p > 0$ であり、すべての元 $\alpha \in k'$ に対して
$\alpha^q \in k$ となる $p$ の冪 $q$ が存在するとき、代数拡大
$k'/k$ は{\it 純非分離的}であるという。
\item すべての $\alpha \in k'$ に対して、$k$ 上の $\alpha$ の最小多項式
$P(T) \in k[T]$ が $k'$ 上一次因子に完全に分解するとき、代数拡大
$k'/k$ は{\it 正規}であるという。
\item 代数拡大 $k'/k$ が分離的かつ正規であるとき、{\it ガロア}であるという。
\end{enumerate}
\end{definition}

\noindent
次の補題は、ほかに適切な置き場所がないように思われる。

\begin{lemma}
\label{fields-lemma-pth-root}
$K$ を標数 $p > 0$ の体とする。$L/K$ を分離代数拡大とし、
$\alpha \in L$ とする。
\begin{enumerate}
\item $K$ 上の $\alpha$ の最小多項式の係数が $K$ における $p$ 乗ならば、
$\alpha$ は $L$ における $p$ 乗である。
\item より一般に、$P \in K[T]$ が、(a) $\alpha$ は $P$ の根である、
(b) $P$ は代数閉包内で相異なる根をもつ、(c) $P$ のすべての係数は
$p$ 乗である、という条件を満たす多項式ならば、$\alpha$ は $L$ における
$p$ 乗である。
\end{enumerate}
\end{lemma}

\begin{proof}
定義から (2) が (1) を導く。$P$ は (2) のとおりであると仮定する。
$P(T) = \sum\nolimits_{i = 0}^d a_i T^{d - i}$ および $a_i = b_i^p$ と書く。
多項式 $Q(T) = \sum\nolimits_{i = 0}^d b_i T^{d - i}$ も代数閉包内で
相異なる根をもつ。$Q$ の根は $P$ の根の $p$ 乗根だからである。
$\alpha$ が $p$ 乗でなければ、$T^p - \alpha$ は $L$ 上の既約多項式である
（補題 \ref{fields-lemma-take-pth-root}）。さらに $Q$ と $T^p - \alpha$ は
代数閉包 $\overline{L}$ 内に共通根をもつ。したがって $Q$ と
$T^p - \alpha$ は互いに素ではなく、$L[T]$ において
$T^p - \alpha | Q$ である。これは $Q$ の根が相異なることに矛盾する。
\end{proof}
